\section{Learning Objectives}

After completing this chapter, readers will be able to:
\begin{itemize}
    \item Implement comprehensive model versioning strategies with MLflow
    \item Track model lineage, metadata, and dependencies systematically
    \item Design and implement approval workflows for model governance
    \item Compare models across multiple dimensions and promote winners
    \item Integrate model registry with CI/CD pipelines for automation
    \item Package and containerize models for reproducible deployment
    \item Build production-grade model management systems
\end{itemize}

\section{Introduction}

Model versioning and registry systems are critical infrastructure for production ML. As teams deploy dozens or hundreds of model versions, managing which models are in production, tracking their lineage, and governing their lifecycle becomes essential. Without proper model management, teams face chaos: lost experiments, unclear model provenance, ungoverned deployments, and inability to reproduce or rollback models.

A robust model registry provides:
\begin{itemize}
    \item \textbf{Version control:} Track every model version with immutable identifiers
    \item \textbf{Lineage tracking:} Understand how models were trained, from data to hyperparameters
    \item \textbf{Governance:} Approval workflows ensure only validated models reach production
    \item \textbf{Reproducibility:} Package models with dependencies for exact recreation
    \item \textbf{Comparison:} Evaluate models systematically before promotion
    \item \textbf{Audit trail:} Complete history of model changes for compliance
\end{itemize}

This chapter covers production-grade model management systems, centered on MLflow but applicable to any registry system.

\section{MLflow Model Registry: Advanced Features}

\subsection{Architecture and Core Concepts}

MLflow Model Registry provides centralized model storage with versioning, stage transitions, and annotations. Understanding its architecture enables effective model management at scale.

\textbf{Key concepts:}
\begin{itemize}
    \item \textbf{Registered Model:} A logical model entity with multiple versions
    \item \textbf{Model Version:} A specific trained instance with unique version number
    \item \textbf{Model Stage:} Lifecycle stage (None, Staging, Production, Archived)
    \item \textbf{Model Signature:} Input/output schema for validation
    \item \textbf{Model Artifact:} Serialized model files and dependencies
\end{itemize}

\subsection{Complete MLflow Registry Implementation}

\begin{lstlisting}[language=Python, caption=Advanced MLflow Model Registry Manager]
import mlflow
from mlflow.tracking import MlflowClient
from mlflow.models.signature import infer_signature
from typing import Dict, List, Optional, Any
from dataclasses import dataclass
from datetime import datetime
import pandas as pd
import numpy as np
import logging
import json

@dataclass
class ModelMetadata:
    """Comprehensive model metadata"""
    name: str
    version: str
    stage: str
    run_id: str
    experiment_id: str

    # Performance metrics
    metrics: Dict[str, float]

    # Training metadata
    training_date: datetime
    training_duration_seconds: float
    dataset_version: str
    dataset_size: int

    # Model characteristics
    model_type: str
    framework: str
    framework_version: str

    # Environment
    python_version: str
    dependencies: List[str]

    # Tags and annotations
    tags: Dict[str, str]
    description: str

class MLflowModelRegistry:
    """
    Advanced MLflow Model Registry manager

    Provides high-level API for model management, versioning, and governance
    """

    def __init__(self, tracking_uri: str = None, registry_uri: str = None):
        """
        Initialize MLflow registry manager

        Args:
            tracking_uri: MLflow tracking server URI
            registry_uri: Model registry URI (defaults to tracking_uri)
        """
        if tracking_uri:
            mlflow.set_tracking_uri(tracking_uri)
        if registry_uri:
            mlflow.set_registry_uri(registry_uri)

        self.client = MlflowClient()
        self.logger = logging.getLogger("MLflowRegistry")

    def register_model(self,
                       model,
                       model_name: str,
                       run_id: str,
                       metadata: ModelMetadata,
                       input_example: Any = None,
                       signature: Any = None) -> str:
        """
        Register a new model version with comprehensive metadata

        Args:
            model: Trained model object
            model_name: Registered model name
            run_id: MLflow run ID
            metadata: Model metadata
            input_example: Example input for schema inference
            signature: Model signature (inferred if not provided)

        Returns:
            Model version number
        """
        self.logger.info(f"Registering model: {model_name}")

        # Infer signature if not provided
        if signature is None and input_example is not None:
            # Generate predictions for signature
            predictions = model.predict(input_example)
            signature = infer_signature(input_example, predictions)

        # Log model with MLflow
        with mlflow.start_run(run_id=run_id):
            # Log model
            mlflow.sklearn.log_model(
                model,
                artifact_path="model",
                registered_model_name=model_name,
                signature=signature,
                input_example=input_example
            )

            # Log metrics
            for metric_name, metric_value in metadata.metrics.items():
                mlflow.log_metric(metric_name, metric_value)

            # Log parameters
            mlflow.log_param("model_type", metadata.model_type)
            mlflow.log_param("framework", metadata.framework)
            mlflow.log_param("framework_version", metadata.framework_version)
            mlflow.log_param("dataset_version", metadata.dataset_version)
            mlflow.log_param("dataset_size", metadata.dataset_size)
            mlflow.log_param("training_duration", metadata.training_duration_seconds)

            # Log tags
            for tag_key, tag_value in metadata.tags.items():
                mlflow.set_tag(tag_key, tag_value)

        # Get the latest version
        latest_versions = self.client.get_latest_versions(model_name)
        if latest_versions:
            version = latest_versions[0].version

            # Add description
            self.client.update_model_version(
                name=model_name,
                version=version,
                description=metadata.description
            )

            self.logger.info(f"Registered {model_name} version {version}")
            return version

        return None

    def transition_model_stage(self,
                              model_name: str,
                              version: str,
                              stage: str,
                              archive_existing: bool = True) -> bool:
        """
        Transition model version to new stage

        Args:
            model_name: Registered model name
            version: Model version
            stage: Target stage (Staging, Production, Archived)
            archive_existing: Archive existing models in target stage

        Returns:
            True if transition successful
        """
        try:
            self.client.transition_model_version_stage(
                name=model_name,
                version=version,
                stage=stage,
                archive_existing_versions=archive_existing
            )

            self.logger.info(
                f"Transitioned {model_name} v{version} to {stage}"
            )
            return True

        except Exception as e:
            self.logger.error(f"Stage transition failed: {str(e)}")
            return False

    def get_model_version(self,
                         model_name: str,
                         version: str = None,
                         stage: str = None) -> Optional[Dict]:
        """
        Get model version details

        Args:
            model_name: Registered model name
            version: Specific version (optional)
            stage: Get latest version in stage (optional)

        Returns:
            Model version information
        """
        try:
            if version:
                mv = self.client.get_model_version(model_name, version)
            elif stage:
                versions = self.client.get_latest_versions(model_name, stages=[stage])
                if not versions:
                    return None
                mv = versions[0]
            else:
                return None

            # Get run details
            run = self.client.get_run(mv.run_id)

            return {
                'name': mv.name,
                'version': mv.version,
                'stage': mv.current_stage,
                'run_id': mv.run_id,
                'description': mv.description,
                'tags': mv.tags,
                'metrics': run.data.metrics,
                'params': run.data.params,
                'creation_timestamp': mv.creation_timestamp,
                'last_updated_timestamp': mv.last_updated_timestamp,
                'source': mv.source
            }

        except Exception as e:
            self.logger.error(f"Failed to get model version: {str(e)}")
            return None

    def compare_models(self,
                      model_name: str,
                      version1: str,
                      version2: str,
                      metric_names: List[str]) -> Dict:
        """
        Compare two model versions across multiple metrics

        Args:
            model_name: Registered model name
            version1: First version to compare
            version2: Second version to compare
            metric_names: Metrics to compare

        Returns:
            Comparison results
        """
        v1_info = self.get_model_version(model_name, version1)
        v2_info = self.get_model_version(model_name, version2)

        if not v1_info or not v2_info:
            return {'error': 'One or both versions not found'}

        comparison = {
            'model_name': model_name,
            'version1': version1,
            'version2': version2,
            'metrics': {}
        }

        for metric_name in metric_names:
            m1 = v1_info['metrics'].get(metric_name)
            m2 = v2_info['metrics'].get(metric_name)

            if m1 is not None and m2 is not None:
                improvement = ((m2 - m1) / abs(m1)) * 100 if m1 != 0 else 0

                comparison['metrics'][metric_name] = {
                    'version1': m1,
                    'version2': m2,
                    'improvement_percent': improvement,
                    'winner': version2 if m2 > m1 else version1
                }

        return comparison

    def search_models(self,
                     filter_string: str = None,
                     max_results: int = 100) -> List[Dict]:
        """
        Search registered models

        Args:
            filter_string: Filter query (e.g., "name LIKE 'fraud%'")
            max_results: Maximum results to return

        Returns:
            List of matching models
        """
        try:
            models = self.client.search_registered_models(
                filter_string=filter_string,
                max_results=max_results
            )

            return [
                {
                    'name': model.name,
                    'creation_timestamp': model.creation_timestamp,
                    'last_updated_timestamp': model.last_updated_timestamp,
                    'description': model.description,
                    'tags': model.tags,
                    'latest_versions': [
                        {
                            'version': v.version,
                            'stage': v.current_stage
                        }
                        for v in model.latest_versions
                    ]
                }
                for model in models
            ]

        except Exception as e:
            self.logger.error(f"Model search failed: {str(e)}")
            return []

    def delete_model_version(self,
                            model_name: str,
                            version: str) -> bool:
        """
        Delete a model version

        Args:
            model_name: Registered model name
            version: Version to delete

        Returns:
            True if deletion successful
        """
        try:
            self.client.delete_model_version(model_name, version)
            self.logger.info(f"Deleted {model_name} version {version}")
            return True

        except Exception as e:
            self.logger.error(f"Deletion failed: {str(e)}")
            return False

    def get_model_stages_summary(self, model_name: str) -> Dict:
        """
        Get summary of all model versions by stage

        Args:
            model_name: Registered model name

        Returns:
            Summary by stage
        """
        try:
            all_versions = self.client.search_model_versions(
                f"name='{model_name}'"
            )

            summary = {
                'Production': [],
                'Staging': [],
                'Archived': [],
                'None': []
            }

            for mv in all_versions:
                summary[mv.current_stage].append({
                    'version': mv.version,
                    'run_id': mv.run_id,
                    'creation_timestamp': mv.creation_timestamp
                })

            return summary

        except Exception as e:
            self.logger.error(f"Failed to get stages summary: {str(e)}")
            return {}

    def load_production_model(self, model_name: str):
        """
        Load the current production model

        Args:
            model_name: Registered model name

        Returns:
            Loaded model
        """
        model_uri = f"models:/{model_name}/Production"
        return mlflow.sklearn.load_model(model_uri)

# Example usage
def example_mlflow_registry():
    """Example of using MLflow Model Registry"""

    # Initialize registry
    registry = MLflowModelRegistry(
        tracking_uri="http://localhost:5000"
    )

    # Train a model (simplified)
    from sklearn.ensemble import RandomForestClassifier
    from sklearn.datasets import make_classification

    X, y = make_classification(n_samples=1000, n_features=20, random_state=42)
    model = RandomForestClassifier(n_estimators=100, random_state=42)
    model.fit(X, y)

    # Create metadata
    metadata = ModelMetadata(
        name="fraud_detector",
        version="1",
        stage="None",
        run_id=mlflow.active_run().info.run_id,
        experiment_id="0",
        metrics={
            'accuracy': 0.95,
            'precision': 0.93,
            'recall': 0.97,
            'f1_score': 0.95
        },
        training_date=datetime.now(),
        training_duration_seconds=120.5,
        dataset_version="v2023-01-15",
        dataset_size=1000,
        model_type="RandomForest",
        framework="scikit-learn",
        framework_version="1.4.2",
        python_version="3.11.6",
        dependencies=["scikit-learn==1.4.2", "numpy==1.26.4"],
        tags={'team': 'ml-team', 'use_case': 'fraud_detection'},
        description="Random Forest model for fraud detection v1"
    )

    # Register model
    version = registry.register_model(
        model=model,
        model_name="fraud_detector",
        run_id=mlflow.active_run().info.run_id,
        metadata=metadata,
        input_example=X[:5]
    )

    print(f"Registered model version: {version}")

    # Transition to staging
    registry.transition_model_stage(
        model_name="fraud_detector",
        version=version,
        stage="Staging"
    )

    # Get model info
    model_info = registry.get_model_version(
        model_name="fraud_detector",
        stage="Staging"
    )
    print(f"Staging model: {model_info}")

    # Compare models (if multiple versions exist)
    # comparison = registry.compare_models(
    #     model_name="fraud_detector",
    #     version1="1",
    #     version2="2",
    #     metric_names=['accuracy', 'f1_score']
    # )

    # Get stages summary
    summary = registry.get_model_stages_summary("fraud_detector")
    print(f"Model stages: {summary}")
\end{lstlisting}

\subsection{Model Signature and Schema Validation}

Model signatures enforce input/output schemas, preventing silent failures from schema mismatches.

\begin{lstlisting}[language=Python, caption=Model Signature and Validation]
from typing import Any, List
from mlflow.models.signature import ModelSignature, infer_signature
from mlflow.types.schema import Schema, ColSpec
import pandas as pd
import numpy as np
from sklearn.ensemble import RandomForestClassifier

class ModelSignatureValidator:
    """Validate model inputs against registered signature"""

    @staticmethod
    def create_signature(input_df: pd.DataFrame,
                        output: Any) -> ModelSignature:
        """
        Create explicit model signature

        Args:
            input_df: Example input DataFrame
            output: Example output

        Returns:
            ModelSignature
        """
        return infer_signature(input_df, output)

    @staticmethod
    def create_manual_signature(
        input_schema: List[tuple],
        output_schema: List[tuple]
    ) -> ModelSignature:
        """
        Create signature manually

        Args:
            input_schema: List of (name, type) tuples
            output_schema: List of (name, type) tuples

        Returns:
            ModelSignature
        """
        input_cols = [ColSpec(dtype, name) for name, dtype in input_schema]
        output_cols = [ColSpec(dtype, name) for name, dtype in output_schema]

        return ModelSignature(
            inputs=Schema(input_cols),
            outputs=Schema(output_cols)
        )

    @staticmethod
    def validate_input(input_df: pd.DataFrame,
                      signature: ModelSignature) -> tuple:
        """
        Validate input against signature

        Args:
            input_df: Input DataFrame
            signature: Expected signature

        Returns:
            (is_valid, error_message)
        """
        if signature is None or signature.inputs is None:
            return True, None

        # Check columns
        expected_cols = set(signature.inputs.input_names())
        actual_cols = set(input_df.columns)

        missing_cols = expected_cols - actual_cols
        extra_cols = actual_cols - expected_cols

        if missing_cols:
            return False, f"Missing columns: {missing_cols}"

        if extra_cols:
            return False, f"Unexpected columns: {extra_cols}"

        # Check types
        for col_spec in signature.inputs.inputs:
            if col_spec.name in input_df.columns:
                actual_dtype = input_df[col_spec.name].dtype
                # Type checking logic here
                pass

        return True, None

# Example usage
def example_signature_validation():
    """Example of signature validation"""

    # Create sample data
    X_train = pd.DataFrame({
        'feature_1': np.random.rand(100),
        'feature_2': np.random.rand(100),
        'feature_3': np.random.randint(0, 10, 100)
    })

    y_train = np.random.randint(0, 2, 100)

    # Train model
    model = RandomForestClassifier()
    model.fit(X_train, y_train)

    # Create signature
    predictions = model.predict(X_train[:5])
    signature = infer_signature(X_train, predictions)

    print(f"Signature: {signature}")

    # Validate input
    validator = ModelSignatureValidator()

    # Valid input
    valid_input = pd.DataFrame({
        'feature_1': [0.5],
        'feature_2': [0.3],
        'feature_3': [5]
    })

    is_valid, error = validator.validate_input(valid_input, signature)
    print(f"Valid input: {is_valid}")

    # Invalid input (missing column)
    invalid_input = pd.DataFrame({
        'feature_1': [0.5],
        'feature_2': [0.3]
    })

    is_valid, error = validator.validate_input(invalid_input, signature)
    print(f"Invalid input: {is_valid}, Error: {error}")
\end{lstlisting}

\section{Model Lineage and Metadata Tracking}

\subsection{Comprehensive Lineage System}

Model lineage tracks the complete history of how a model was created, enabling reproducibility and debugging.

\begin{lstlisting}[language=Python, caption=Model Lineage Tracker]
from typing import Any, Dict, List, Optional, Set
from dataclasses import dataclass, field
from datetime import datetime
import hashlib
import json
import logging
import numpy as np
import pandas as pd
import mlflow
import mlflow

@dataclass
class DataLineage:
    """Data lineage information"""
    dataset_name: str
    dataset_version: str
    dataset_hash: str
    num_samples: int
    num_features: int
    feature_names: List[str]
    data_source: str
    collection_date: datetime
    preprocessing_steps: List[str] = field(default_factory=list)

@dataclass
class CodeLineage:
    """Code lineage information"""
    git_commit: str
    git_branch: str
    git_repo: str
    code_version: str
    file_paths: List[str]
    dependencies: Dict[str, str]

@dataclass
class ModelLineage:
    """Complete model lineage"""
    model_id: str
    model_name: str
    model_version: str
    creation_timestamp: datetime

    # Data lineage
    data_lineage: DataLineage

    # Code lineage
    code_lineage: CodeLineage

    # Training configuration
    hyperparameters: Dict[str, Any]
    training_duration: float
    hardware: Dict[str, str]

    # Parent models (for transfer learning, ensembles)
    parent_models: List[str] = field(default_factory=list)

    # Performance
    metrics: Dict[str, float] = field(default_factory=dict)

    # Additional metadata
    tags: Dict[str, str] = field(default_factory=dict)
    notes: str = ""

class LineageTracker:
    """
    Track comprehensive model lineage

    Captures data, code, and training provenance
    """

    def __init__(self, storage_backend: str = "mlflow"):
        """
        Initialize lineage tracker

        Args:
            storage_backend: Where to store lineage (mlflow, database, etc.)
        """
        self.storage_backend = storage_backend
        self.logger = logging.getLogger("LineageTracker")

    def track_data_lineage(self,
                          dataset: pd.DataFrame,
                          dataset_name: str,
                          dataset_version: str,
                          data_source: str,
                          preprocessing_steps: List[str] = None) -> DataLineage:
        """
        Track data lineage

        Args:
            dataset: Training dataset
            dataset_name: Dataset identifier
            dataset_version: Dataset version
            data_source: Where data came from
            preprocessing_steps: List of preprocessing operations

        Returns:
            DataLineage object
        """
        # Calculate dataset hash for verification
        dataset_hash = hashlib.sha256(
            pd.util.hash_pandas_object(dataset).values
        ).hexdigest()

        return DataLineage(
            dataset_name=dataset_name,
            dataset_version=dataset_version,
            dataset_hash=dataset_hash,
            num_samples=len(dataset),
            num_features=len(dataset.columns),
            feature_names=list(dataset.columns),
            data_source=data_source,
            collection_date=datetime.now(),
            preprocessing_steps=preprocessing_steps or []
        )

    def track_code_lineage(self,
                          training_script_paths: List[str]) -> CodeLineage:
        """
        Track code lineage from git

        Args:
            training_script_paths: Paths to training scripts

        Returns:
            CodeLineage object
        """
        import subprocess
        import pkg_resources

        try:
            # Get git information
            git_commit = subprocess.check_output(
                ['git', 'rev-parse', 'HEAD']
            ).decode('utf-8').strip()

            git_branch = subprocess.check_output(
                ['git', 'rev-parse', '--abbrev-ref', 'HEAD']
            ).decode('utf-8').strip()

            git_repo = subprocess.check_output(
                ['git', 'config', '--get', 'remote.origin.url']
            ).decode('utf-8').strip()

        except Exception as e:
            self.logger.warning(f"Failed to get git info: {str(e)}")
            git_commit = "unknown"
            git_branch = "unknown"
            git_repo = "unknown"

        # Get installed packages
        dependencies = {
            pkg.key: pkg.version
            for pkg in pkg_resources.working_set
        }

        return CodeLineage(
            git_commit=git_commit,
            git_branch=git_branch,
            git_repo=git_repo,
            code_version="1.0",
            file_paths=training_script_paths,
            dependencies=dependencies
        )

    def create_lineage(self,
                      model_id: str,
                      model_name: str,
                      model_version: str,
                      data_lineage: DataLineage,
                      code_lineage: CodeLineage,
                      hyperparameters: Dict,
                      training_duration: float,
                      metrics: Dict[str, float],
                      parent_models: List[str] = None) -> ModelLineage:
        """
        Create complete model lineage

        Args:
            model_id: Unique model identifier
            model_name: Model name
            model_version: Model version
            data_lineage: Data lineage
            code_lineage: Code lineage
            hyperparameters: Training hyperparameters
            training_duration: Training time in seconds
            metrics: Performance metrics
            parent_models: Parent model IDs (for transfer learning)

        Returns:
            Complete ModelLineage
        """
        import platform

        lineage = ModelLineage(
            model_id=model_id,
            model_name=model_name,
            model_version=model_version,
            creation_timestamp=datetime.now(),
            data_lineage=data_lineage,
            code_lineage=code_lineage,
            hyperparameters=hyperparameters,
            training_duration=training_duration,
            hardware={
                'platform': platform.system(),
                'processor': platform.processor(),
                'python_version': platform.python_version()
            },
            parent_models=parent_models or [],
            metrics=metrics
        )

        # Store lineage
        self._store_lineage(lineage)

        return lineage

    def _store_lineage(self, lineage: ModelLineage):
        """Store lineage in backend"""
        if self.storage_backend == "mlflow":
            self._store_in_mlflow(lineage)
        else:
            self._store_in_json(lineage)

    def _store_in_mlflow(self, lineage: ModelLineage):
        """Store lineage in MLflow"""
        with mlflow.start_run(run_name=f"{lineage.model_name}_v{lineage.model_version}"):
            # Log data lineage
            mlflow.log_param("dataset_name", lineage.data_lineage.dataset_name)
            mlflow.log_param("dataset_version", lineage.data_lineage.dataset_version)
            mlflow.log_param("dataset_hash", lineage.data_lineage.dataset_hash)
            mlflow.log_param("num_samples", lineage.data_lineage.num_samples)
            mlflow.log_param("num_features", lineage.data_lineage.num_features)

            # Log code lineage
            mlflow.log_param("git_commit", lineage.code_lineage.git_commit)
            mlflow.log_param("git_branch", lineage.code_lineage.git_branch)

            # Log hyperparameters
            for key, value in lineage.hyperparameters.items():
                mlflow.log_param(key, value)

            # Log metrics
            for key, value in lineage.metrics.items():
                mlflow.log_metric(key, value)

            # Log lineage as artifact
            lineage_dict = self._lineage_to_dict(lineage)
            with open("lineage.json", "w") as f:
                json.dump(lineage_dict, f, indent=2, default=str)
            mlflow.log_artifact("lineage.json")

    def _store_in_json(self, lineage: ModelLineage):
        """Store lineage as JSON file"""
        lineage_dict = self._lineage_to_dict(lineage)
        filename = f"lineage_{lineage.model_id}.json"
        with open(filename, "w") as f:
            json.dump(lineage_dict, f, indent=2, default=str)

    def _lineage_to_dict(self, lineage: ModelLineage) -> Dict:
        """Convert lineage to dictionary"""
        return {
            'model_id': lineage.model_id,
            'model_name': lineage.model_name,
            'model_version': lineage.model_version,
            'creation_timestamp': lineage.creation_timestamp.isoformat(),
            'data_lineage': {
                'dataset_name': lineage.data_lineage.dataset_name,
                'dataset_version': lineage.data_lineage.dataset_version,
                'dataset_hash': lineage.data_lineage.dataset_hash,
                'num_samples': lineage.data_lineage.num_samples,
                'num_features': lineage.data_lineage.num_features,
                'feature_names': lineage.data_lineage.feature_names,
                'data_source': lineage.data_lineage.data_source,
                'preprocessing_steps': lineage.data_lineage.preprocessing_steps
            },
            'code_lineage': {
                'git_commit': lineage.code_lineage.git_commit,
                'git_branch': lineage.code_lineage.git_branch,
                'git_repo': lineage.code_lineage.git_repo,
                'code_version': lineage.code_lineage.code_version,
                'file_paths': lineage.code_lineage.file_paths
            },
            'hyperparameters': lineage.hyperparameters,
            'training_duration': lineage.training_duration,
            'hardware': lineage.hardware,
            'parent_models': lineage.parent_models,
            'metrics': lineage.metrics,
            'tags': lineage.tags,
            'notes': lineage.notes
        }

    def get_lineage(self, model_id: str) -> Optional[ModelLineage]:
        """
        Retrieve model lineage

        Args:
            model_id: Model identifier

        Returns:
            ModelLineage or None
        """
        try:
            with open(f"lineage_{model_id}.json", "r") as f:
                lineage_dict = json.load(f)
                # Reconstruct ModelLineage from dict
                # (Simplified - full implementation would deserialize properly)
                return lineage_dict
        except FileNotFoundError:
            return None

# Example usage
def example_lineage_tracking():
    """Example of comprehensive lineage tracking"""

    # Initialize tracker
    tracker = LineageTracker(storage_backend="mlflow")

    # Create training dataset
    X_train = pd.DataFrame({
        'age': np.random.randint(18, 80, 1000),
        'income': np.random.randint(20000, 200000, 1000),
        'credit_score': np.random.randint(300, 850, 1000)
    })

    # Track data lineage
    data_lineage = tracker.track_data_lineage(
        dataset=X_train,
        dataset_name="customer_features",
        dataset_version="v2024-01-15",
        data_source="data_warehouse.customers",
        preprocessing_steps=[
            "remove_nulls",
            "normalize_income",
            "handle_outliers"
        ]
    )

    # Track code lineage
    code_lineage = tracker.track_code_lineage(
        training_script_paths=["train.py", "preprocess.py"]
    )

    # Training configuration
    hyperparameters = {
        'n_estimators': 100,
        'max_depth': 10,
        'min_samples_split': 5,
        'learning_rate': 0.1
    }

    # Create complete lineage
    lineage = tracker.create_lineage(
        model_id="model_abc123",
        model_name="credit_risk_model",
        model_version="v2.0",
        data_lineage=data_lineage,
        code_lineage=code_lineage,
        hyperparameters=hyperparameters,
        training_duration=450.5,
        metrics={
            'accuracy': 0.92,
            'auc': 0.88,
            'f1_score': 0.90
        },
        parent_models=[]
    )

    print(f"Lineage created for {lineage.model_id}")
    print(f"Dataset: {lineage.data_lineage.dataset_name}")
    print(f"Git commit: {lineage.code_lineage.git_commit}")
\end{lstlisting}

\section{Approval Workflows and Governance}

\subsection{Model Governance Framework}

Production ML systems require governance to ensure only validated models reach production. Approval workflows enforce reviews, validations, and compliance checks before model promotion.

\textbf{Key governance requirements:}
\begin{itemize}
    \item \textbf{Multi-stage approvals:} Different stakeholders approve at each stage
    \item \textbf{Validation gates:} Automated checks before promotion
    \item \textbf{Role-based access:} Control who can promote models
    \item \textbf{Audit trails:} Complete history of approvals and rejections
    \item \textbf{Compliance integration:} Ticketing systems for regulatory requirements
    \item \textbf{Automated enforcement:} Prevent unauthorized promotions
\end{itemize}

\subsection{Complete Governance Implementation}

\begin{lstlisting}[language=Python, caption=Model Approval and Governance System]
from enum import Enum
from typing import Any, Dict, List, Optional, Callable
from dataclasses import dataclass, field
from datetime import datetime
import logging
import json

class ApprovalStatus(Enum):
    """Approval request status"""
    PENDING = "pending"
    APPROVED = "approved"
    REJECTED = "rejected"
    CANCELLED = "cancelled"

class UserRole(Enum):
    """User roles for RBAC"""
    DATA_SCIENTIST = "data_scientist"
    ML_ENGINEER = "ml_engineer"
    TEAM_LEAD = "team_lead"
    COMPLIANCE_OFFICER = "compliance_officer"
    PRODUCTION_OWNER = "production_owner"

@dataclass
class ValidationGate:
    """Validation gate configuration"""
    name: str
    description: str
    validator_func: Callable
    required: bool = True
    error_message: str = ""

@dataclass
class ApprovalRequest:
    """Model approval request"""
    request_id: str
    model_name: str
    current_version: str
    source_stage: str
    target_stage: str

    # Requester information
    requester: str
    requester_role: UserRole
    request_timestamp: datetime

    # Approval chain
    required_approvers: List[UserRole]
    approvals: List[Dict[str, Any]] = field(default_factory=list)
    rejections: List[Dict[str, Any]] = field(default_factory=list)

    # Status
    status: ApprovalStatus = ApprovalStatus.PENDING

    # Validation results
    validation_results: Dict[str, bool] = field(default_factory=dict)

    # Metadata
    justification: str = ""
    impact_assessment: str = ""
    rollback_plan: str = ""
    ticket_id: Optional[str] = None

class ModelGovernanceManager:
    """
    Comprehensive model governance and approval workflow manager

    Enforces approval workflows, validation gates, and RBAC
    """

    def __init__(self,
                 mlflow_registry: MLflowModelRegistry,
                 ticket_system: Optional[Any] = None):
        """
        Initialize governance manager

        Args:
            mlflow_registry: MLflow registry instance
            ticket_system: Optional ticketing system integration
        """
        self.registry = mlflow_registry
        self.ticket_system = ticket_system
        self.logger = logging.getLogger("ModelGovernance")

        # Validation gates by stage transition
        self.validation_gates: Dict[str, List[ValidationGate]] = {}

        # Approval requirements by stage
        self.approval_requirements: Dict[str, List[UserRole]] = {
            'Staging': [UserRole.TEAM_LEAD],
            'Production': [UserRole.TEAM_LEAD, UserRole.PRODUCTION_OWNER],
            'Archived': [UserRole.TEAM_LEAD]
        }

        # Approval requests
        self.approval_requests: Dict[str, ApprovalRequest] = {}

        # Audit log
        self.audit_log: List[Dict] = []

    def register_validation_gate(self,
                                 target_stage: str,
                                 gate: ValidationGate):
        """
        Register validation gate for stage transition

        Args:
            target_stage: Target stage (Staging, Production, etc.)
            gate: ValidationGate to register
        """
        if target_stage not in self.validation_gates:
            self.validation_gates[target_stage] = []

        self.validation_gates[target_stage].append(gate)
        self.logger.info(f"Registered validation gate '{gate.name}' for {target_stage}")

    def request_approval(self,
                        model_name: str,
                        version: str,
                        source_stage: str,
                        target_stage: str,
                        requester: str,
                        requester_role: UserRole,
                        justification: str = "",
                        impact_assessment: str = "",
                        rollback_plan: str = "") -> str:
        """
        Request approval for model stage transition

        Args:
            model_name: Model name
            version: Model version
            source_stage: Current stage
            target_stage: Target stage
            requester: User requesting approval
            requester_role: Requester's role
            justification: Reason for promotion
            impact_assessment: Expected impact
            rollback_plan: Plan if rollback needed

        Returns:
            Request ID
        """
        import uuid

        request_id = str(uuid.uuid4())

        # Get required approvers for target stage
        required_approvers = self.approval_requirements.get(
            target_stage,
            [UserRole.TEAM_LEAD]
        )

        # Create approval request
        request = ApprovalRequest(
            request_id=request_id,
            model_name=model_name,
            current_version=version,
            source_stage=source_stage,
            target_stage=target_stage,
            requester=requester,
            requester_role=requester_role,
            request_timestamp=datetime.now(),
            required_approvers=required_approvers,
            justification=justification,
            impact_assessment=impact_assessment,
            rollback_plan=rollback_plan
        )

        # Run validation gates
        validation_passed = self._run_validation_gates(
            request,
            model_name,
            version
        )

        if not validation_passed:
            request.status = ApprovalStatus.REJECTED
            self.logger.warning(
                f"Validation failed for {model_name} v{version}"
            )
            self.approval_requests[request_id] = request
            return request_id

        # Create ticket if system available
        if self.ticket_system:
            ticket_id = self._create_governance_ticket(request)
            request.ticket_id = ticket_id

        self.approval_requests[request_id] = request

        # Log audit event
        self._log_audit_event(
            event_type="approval_requested",
            request_id=request_id,
            model_name=model_name,
            version=version,
            user=requester,
            details={
                'source_stage': source_stage,
                'target_stage': target_stage,
                'justification': justification
            }
        )

        self.logger.info(
            f"Approval request created: {request_id} for {model_name} v{version}"
        )

        return request_id

    def approve_request(self,
                       request_id: str,
                       approver: str,
                       approver_role: UserRole,
                       comments: str = "") -> bool:
        """
        Approve a model promotion request

        Args:
            request_id: Request ID
            approver: User approving
            approver_role: Approver's role
            comments: Approval comments

        Returns:
            True if approval recorded successfully
        """
        request = self.approval_requests.get(request_id)
        if not request:
            self.logger.error(f"Request {request_id} not found")
            return False

        if request.status != ApprovalStatus.PENDING:
            self.logger.error(
                f"Request {request_id} is not pending (status: {request.status})"
            )
            return False

        # Check if approver role is required
        if approver_role not in request.required_approvers:
            self.logger.error(
                f"{approver_role} not authorized to approve this request"
            )
            return False

        # Record approval
        approval = {
            'approver': approver,
            'role': approver_role.value,
            'timestamp': datetime.now().isoformat(),
            'comments': comments
        }
        request.approvals.append(approval)

        # Log audit event
        self._log_audit_event(
            event_type="approval_granted",
            request_id=request_id,
            model_name=request.model_name,
            version=request.current_version,
            user=approver,
            details={'role': approver_role.value, 'comments': comments}
        )

        # Check if all required approvals received
        approved_roles = {a['role'] for a in request.approvals}
        required_roles = {r.value for r in request.required_approvers}

        if required_roles.issubset(approved_roles):
            # All approvals received - execute promotion
            success = self._execute_promotion(request)
            if success:
                request.status = ApprovalStatus.APPROVED
                self.logger.info(
                    f"Request {request_id} fully approved and executed"
                )
            return success

        self.logger.info(
            f"Approval recorded for {request_id}. "
            f"Pending: {required_roles - approved_roles}"
        )
        return True

    def reject_request(self,
                      request_id: str,
                      rejector: str,
                      rejector_role: UserRole,
                      reason: str) -> bool:
        """
        Reject a model promotion request

        Args:
            request_id: Request ID
            rejector: User rejecting
            rejector_role: Rejector's role
            reason: Rejection reason

        Returns:
            True if rejection recorded successfully
        """
        request = self.approval_requests.get(request_id)
        if not request:
            self.logger.error(f"Request {request_id} not found")
            return False

        if request.status != ApprovalStatus.PENDING:
            self.logger.error(
                f"Request {request_id} is not pending"
            )
            return False

        # Record rejection
        rejection = {
            'rejector': rejector,
            'role': rejector_role.value,
            'timestamp': datetime.now().isoformat(),
            'reason': reason
        }
        request.rejections.append(rejection)
        request.status = ApprovalStatus.REJECTED

        # Log audit event
        self._log_audit_event(
            event_type="approval_rejected",
            request_id=request_id,
            model_name=request.model_name,
            version=request.current_version,
            user=rejector,
            details={'role': rejector_role.value, 'reason': reason}
        )

        # Update ticket if exists
        if self.ticket_system and request.ticket_id:
            self._update_ticket_status(request.ticket_id, "Rejected", reason)

        self.logger.info(f"Request {request_id} rejected by {rejector}")
        return True

    def _run_validation_gates(self,
                             request: ApprovalRequest,
                             model_name: str,
                             version: str) -> bool:
        """
        Run validation gates for stage transition

        Args:
            request: Approval request
            model_name: Model name
            version: Model version

        Returns:
            True if all validations pass
        """
        gates = self.validation_gates.get(request.target_stage, [])

        if not gates:
            self.logger.info(f"No validation gates for {request.target_stage}")
            return True

        all_passed = True

        for gate in gates:
            try:
                # Get model version info
                model_info = self.registry.get_model_version(
                    model_name,
                    version
                )

                # Run validation
                passed = gate.validator_func(model_info, request)
                request.validation_results[gate.name] = passed

                if not passed and gate.required:
                    self.logger.error(
                        f"Required validation gate '{gate.name}' failed"
                    )
                    all_passed = False
                elif not passed:
                    self.logger.warning(
                        f"Optional validation gate '{gate.name}' failed"
                    )
                else:
                    self.logger.info(f"Validation gate '{gate.name}' passed")

            except Exception as e:
                self.logger.error(
                    f"Validation gate '{gate.name}' error: {str(e)}"
                )
                if gate.required:
                    all_passed = False

        return all_passed

    def _execute_promotion(self, request: ApprovalRequest) -> bool:
        """
        Execute model promotion after approvals

        Args:
            request: Approved request

        Returns:
            True if promotion successful
        """
        success = self.registry.transition_model_stage(
            model_name=request.model_name,
            version=request.current_version,
            stage=request.target_stage,
            archive_existing=True
        )

        if success:
            # Log audit event
            self._log_audit_event(
                event_type="model_promoted",
                request_id=request.request_id,
                model_name=request.model_name,
                version=request.current_version,
                user="system",
                details={
                    'source_stage': request.source_stage,
                    'target_stage': request.target_stage
                }
            )

            # Update ticket
            if self.ticket_system and request.ticket_id:
                self._update_ticket_status(
                    request.ticket_id,
                    "Completed",
                    "Model promoted successfully"
                )

        return success

    def _create_governance_ticket(self, request: ApprovalRequest) -> str:
        """Create governance ticket in external system"""
        # Simplified - would integrate with JIRA, ServiceNow, etc.
        ticket = {
            'title': f"Model Promotion: {request.model_name} v{request.current_version}",
            'description': f"""
Model Promotion Request

Model: {request.model_name}
Version: {request.current_version}
Source Stage: {request.source_stage}
Target Stage: {request.target_stage}

Requester: {request.requester}
Justification: {request.justification}

Impact Assessment:
{request.impact_assessment}

Rollback Plan:
{request.rollback_plan}
            """,
            'type': 'Model Governance',
            'priority': 'High' if request.target_stage == 'Production' else 'Medium'
        }

        # In real implementation, call ticket system API
        ticket_id = f"GOV-{request.request_id[:8]}"
        self.logger.info(f"Created ticket {ticket_id}")
        return ticket_id

    def _update_ticket_status(self, ticket_id: str, status: str, comment: str):
        """Update ticket status"""
        # In real implementation, call ticket system API
        self.logger.info(f"Updated ticket {ticket_id}: {status}")

    def _log_audit_event(self,
                        event_type: str,
                        request_id: str,
                        model_name: str,
                        version: str,
                        user: str,
                        details: Dict):
        """Log audit event"""
        event = {
            'timestamp': datetime.now().isoformat(),
            'event_type': event_type,
            'request_id': request_id,
            'model_name': model_name,
            'version': version,
            'user': user,
            'details': details
        }
        self.audit_log.append(event)

    def get_audit_trail(self,
                       model_name: str = None,
                       start_date: datetime = None,
                       end_date: datetime = None) -> List[Dict]:
        """
        Get audit trail

        Args:
            model_name: Filter by model name
            start_date: Filter by start date
            end_date: Filter by end date

        Returns:
            List of audit events
        """
        filtered_events = self.audit_log

        if model_name:
            filtered_events = [
                e for e in filtered_events
                if e['model_name'] == model_name
            ]

        if start_date:
            filtered_events = [
                e for e in filtered_events
                if datetime.fromisoformat(e['timestamp']) >= start_date
            ]

        if end_date:
            filtered_events = [
                e for e in filtered_events
                if datetime.fromisoformat(e['timestamp']) <= end_date
            ]

        return filtered_events

    def get_request_status(self, request_id: str) -> Optional[Dict]:
        """
        Get approval request status

        Args:
            request_id: Request ID

        Returns:
            Request details
        """
        request = self.approval_requests.get(request_id)
        if not request:
            return None

        return {
            'request_id': request.request_id,
            'model_name': request.model_name,
            'version': request.current_version,
            'source_stage': request.source_stage,
            'target_stage': request.target_stage,
            'status': request.status.value,
            'requester': request.requester,
            'request_timestamp': request.request_timestamp.isoformat(),
            'approvals': request.approvals,
            'rejections': request.rejections,
            'validation_results': request.validation_results,
            'ticket_id': request.ticket_id
        }

# Example validation gates
def validate_model_performance(model_info: Dict, request: ApprovalRequest) -> bool:
    """Validate model meets performance thresholds"""
    metrics = model_info.get('metrics', {})

    # Example thresholds for production
    if request.target_stage == 'Production':
        accuracy = metrics.get('accuracy', 0)
        if accuracy < 0.90:
            return False

        f1_score = metrics.get('f1_score', 0)
        if f1_score < 0.85:
            return False

    return True

def validate_model_testing(model_info: Dict, request: ApprovalRequest) -> bool:
    """Validate model has been tested"""
    tags = model_info.get('tags', {})

    # Check for test completion tag
    if request.target_stage == 'Production':
        if tags.get('tested') != 'true':
            return False

        if tags.get('integration_tested') != 'true':
            return False

    return True

def validate_documentation(model_info: Dict, request: ApprovalRequest) -> bool:
    """Validate model documentation"""
    description = model_info.get('description', '')

    if request.target_stage == 'Production':
        # Require minimum documentation
        if len(description) < 100:
            return False

    return True

# Example usage
def example_governance_workflow():
    """Example of complete governance workflow"""

    # Initialize components
    registry = MLflowModelRegistry(tracking_uri="http://localhost:5000")
    governance = ModelGovernanceManager(registry)

    # Register validation gates for Production
    governance.register_validation_gate(
        target_stage='Production',
        gate=ValidationGate(
            name='performance_check',
            description='Validate model performance metrics',
            validator_func=validate_model_performance,
            required=True
        )
    )

    governance.register_validation_gate(
        target_stage='Production',
        gate=ValidationGate(
            name='testing_check',
            description='Validate testing completion',
            validator_func=validate_model_testing,
            required=True
        )
    )

    governance.register_validation_gate(
        target_stage='Production',
        gate=ValidationGate(
            name='documentation_check',
            description='Validate documentation completeness',
            validator_func=validate_documentation,
            required=False
        )
    )

    # Request approval for production promotion
    request_id = governance.request_approval(
        model_name='fraud_detector',
        version='3',
        source_stage='Staging',
        target_stage='Production',
        requester='alice@company.com',
        requester_role=UserRole.DATA_SCIENTIST,
        justification='New model shows 5% improvement in F1 score',
        impact_assessment='Expected to reduce false positives by 15%',
        rollback_plan='Immediate rollback to v2 if accuracy drops below 0.88'
    )

    print(f"Approval request created: {request_id}")

    # Check request status
    status = governance.get_request_status(request_id)
    print(f"Request status: {status['status']}")
    print(f"Validation results: {status['validation_results']}")

    # Team lead approves
    governance.approve_request(
        request_id=request_id,
        approver='bob@company.com',
        approver_role=UserRole.TEAM_LEAD,
        comments='Model performance looks good, approving for production'
    )

    # Production owner approves (final approval)
    governance.approve_request(
        request_id=request_id,
        approver='charlie@company.com',
        approver_role=UserRole.PRODUCTION_OWNER,
        comments='Infrastructure ready, approved for deployment'
    )

    # Check final status
    final_status = governance.get_request_status(request_id)
    print(f"Final status: {final_status['status']}")

    # Get audit trail
    audit_trail = governance.get_audit_trail(model_name='fraud_detector')
    print(f"Audit events: {len(audit_trail)}")
    for event in audit_trail:
        print(f"  {event['timestamp']}: {event['event_type']} by {event['user']}")
\end{lstlisting}

\section{Model Comparison and Promotion}

\subsection{Comprehensive Model Comparison Framework}

Comparing models rigorously before promotion prevents regression. Effective comparison evaluates models across multiple dimensions: performance metrics, statistical significance, computational cost, data requirements, and operational characteristics.

\textbf{Comparison dimensions:}
\begin{itemize}
    \item \textbf{Performance metrics:} Accuracy, precision, recall, F1, AUC-ROC
    \item \textbf{Statistical significance:} Hypothesis tests, confidence intervals
    \item \textbf{Computational cost:} Inference latency, throughput, resource usage
    \item \textbf{Data requirements:} Feature dependencies, data freshness needs
    \item \textbf{Model complexity:} Number of parameters, interpretability
    \item \textbf{Business metrics:} Revenue impact, user satisfaction, operational cost
\end{itemize}

\subsection{Model Comparison and Promotion System}

\begin{lstlisting}[language=Python, caption=Model Comparison and Promotion System]
from typing import Dict, List, Optional, Tuple
from dataclasses import dataclass
from datetime import datetime
import numpy as np
from scipy import stats
import logging

@dataclass
class ComparisonResult:
    """Results of model comparison"""
    model1_name: str
    model1_version: str
    model2_name: str
    model2_version: str

    # Performance comparison
    metric_comparisons: Dict[str, Dict[str, float]]

    # Statistical significance
    statistically_significant: bool
    p_values: Dict[str, float]
    confidence_intervals: Dict[str, Tuple[float, float]]

    # Computational comparison
    latency_comparison: Dict[str, float]
    throughput_comparison: Dict[str, float]
    resource_comparison: Dict[str, float]

    # Overall recommendation
    recommended_model: str
    recommendation_reason: str
    confidence_score: float

@dataclass
class PromotionCriteria:
    """Criteria for automatic model promotion"""
    # Performance requirements
    min_accuracy_improvement: float = 0.01  # 1% minimum improvement
    min_f1_improvement: float = 0.01
    require_statistical_significance: bool = True
    min_confidence_level: float = 0.95

    # Resource constraints
    max_latency_ms: Optional[float] = None
    max_memory_mb: Optional[float] = None
    min_throughput_qps: Optional[float] = None

    # Business requirements
    max_false_positive_rate: Optional[float] = None
    max_false_negative_rate: Optional[float] = None

    # A/B test requirements
    min_sample_size: int = 1000
    min_test_duration_days: int = 7

class ModelComparator:
    """
    Compare models across multiple dimensions

    Performs statistical testing and comprehensive evaluation
    """

    def __init__(self, mlflow_registry: MLflowModelRegistry):
        """
        Initialize model comparator

        Args:
            mlflow_registry: MLflow registry instance
        """
        self.registry = mlflow_registry
        self.logger = logging.getLogger("ModelComparator")

    def compare_models(self,
                      model_name: str,
                      version1: str,
                      version2: str,
                      test_data: pd.DataFrame = None,
                      test_labels: np.ndarray = None,
                      production_metrics: Dict = None) -> ComparisonResult:
        """
        Comprehensive model comparison

        Args:
            model_name: Model name
            version1: First version
            version2: Second version
            test_data: Test dataset for evaluation
            test_labels: True labels for test data
            production_metrics: Optional production performance data

        Returns:
            ComparisonResult
        """
        self.logger.info(f"Comparing {model_name} v{version1} vs v{version2}")

        # Get model information
        model1_info = self.registry.get_model_version(model_name, version1)
        model2_info = self.registry.get_model_version(model_name, version2)

        # Compare metrics
        metric_comparisons = self._compare_metrics(
            model1_info['metrics'],
            model2_info['metrics']
        )

        # Statistical significance testing
        p_values = {}
        confidence_intervals = {}
        statistically_significant = True

        if test_data is not None and test_labels is not None:
            # Load models and evaluate
            model1 = self._load_model(model_name, version1)
            model2 = self._load_model(model_name, version2)

            pred1 = model1.predict_proba(test_data)[:, 1]
            pred2 = model2.predict_proba(test_data)[:, 1]

            # Perform statistical tests
            stat_results = self._perform_statistical_tests(
                pred1, pred2, test_labels
            )
            p_values = stat_results['p_values']
            confidence_intervals = stat_results['confidence_intervals']
            statistically_significant = stat_results['significant']

        # Compare computational performance
        latency_comparison = self._compare_latency(
            model1_info, model2_info
        )

        throughput_comparison = self._compare_throughput(
            model1_info, model2_info
        )

        resource_comparison = self._compare_resources(
            model1_info, model2_info
        )

        # Make recommendation
        recommended_model, reason, confidence = self._make_recommendation(
            version1, version2,
            metric_comparisons,
            statistically_significant,
            latency_comparison,
            resource_comparison
        )

        return ComparisonResult(
            model1_name=model_name,
            model1_version=version1,
            model2_name=model_name,
            model2_version=version2,
            metric_comparisons=metric_comparisons,
            statistically_significant=statistically_significant,
            p_values=p_values,
            confidence_intervals=confidence_intervals,
            latency_comparison=latency_comparison,
            throughput_comparison=throughput_comparison,
            resource_comparison=resource_comparison,
            recommended_model=recommended_model,
            recommendation_reason=reason,
            confidence_score=confidence
        )

    def _compare_metrics(self,
                        metrics1: Dict[str, float],
                        metrics2: Dict[str, float]) -> Dict[str, Dict[str, float]]:
        """Compare performance metrics"""
        comparisons = {}

        all_metrics = set(metrics1.keys()) | set(metrics2.keys())

        for metric in all_metrics:
            m1 = metrics1.get(metric, 0)
            m2 = metrics2.get(metric, 0)

            improvement = ((m2 - m1) / abs(m1) * 100) if m1 != 0 else 0

            comparisons[metric] = {
                'model1': m1,
                'model2': m2,
                'improvement_percent': improvement,
                'absolute_difference': m2 - m1
            }

        return comparisons

    def _perform_statistical_tests(self,
                                   pred1: np.ndarray,
                                   pred2: np.ndarray,
                                   labels: np.ndarray) -> Dict:
        """Perform statistical significance tests"""

        # Calculate AUC for both models
        from sklearn.metrics import roc_auc_score

        auc1 = roc_auc_score(labels, pred1)
        auc2 = roc_auc_score(labels, pred2)

        # McNemar's test for binary predictions
        binary_pred1 = (pred1 > 0.5).astype(int)
        binary_pred2 = (pred2 > 0.5).astype(int)

        # Create contingency table
        both_correct = np.sum((binary_pred1 == labels) & (binary_pred2 == labels))
        only_1_correct = np.sum((binary_pred1 == labels) & (binary_pred2 != labels))
        only_2_correct = np.sum((binary_pred1 != labels) & (binary_pred2 == labels))
        both_wrong = np.sum((binary_pred1 != labels) & (binary_pred2 != labels))

        # McNemar's test
        if only_1_correct + only_2_correct > 0:
            mcnemar_stat = ((abs(only_1_correct - only_2_correct) - 1) ** 2) / \
                          (only_1_correct + only_2_correct)
            mcnemar_p = 1 - stats.chi2.cdf(mcnemar_stat, 1)
        else:
            mcnemar_p = 1.0

        # Bootstrap confidence interval for AUC difference
        n_bootstrap = 1000
        auc_diffs = []

        for _ in range(n_bootstrap):
            indices = np.random.choice(len(labels), len(labels), replace=True)
            boot_auc1 = roc_auc_score(labels[indices], pred1[indices])
            boot_auc2 = roc_auc_score(labels[indices], pred2[indices])
            auc_diffs.append(boot_auc2 - boot_auc1)

        auc_diffs = np.array(auc_diffs)
        ci_lower = np.percentile(auc_diffs, 2.5)
        ci_upper = np.percentile(auc_diffs, 97.5)

        # Determine if statistically significant
        significant = mcnemar_p < 0.05 and ci_lower > 0

        return {
            'p_values': {
                'mcnemar': mcnemar_p
            },
            'confidence_intervals': {
                'auc_difference': (ci_lower, ci_upper)
            },
            'significant': significant
        }

    def _compare_latency(self,
                        model1_info: Dict,
                        model2_info: Dict) -> Dict[str, float]:
        """Compare inference latency"""
        # In practice, would get from production metrics or benchmarks
        latency1 = model1_info.get('params', {}).get('avg_latency_ms', 0)
        latency2 = model2_info.get('params', {}).get('avg_latency_ms', 0)

        return {
            'model1_ms': float(latency1) if latency1 else 0,
            'model2_ms': float(latency2) if latency2 else 0,
            'difference_ms': float(latency2) - float(latency1) if latency1 and latency2 else 0
        }

    def _compare_throughput(self,
                           model1_info: Dict,
                           model2_info: Dict) -> Dict[str, float]:
        """Compare throughput (queries per second)"""
        throughput1 = model1_info.get('params', {}).get('throughput_qps', 0)
        throughput2 = model2_info.get('params', {}).get('throughput_qps', 0)

        return {
            'model1_qps': float(throughput1) if throughput1 else 0,
            'model2_qps': float(throughput2) if throughput2 else 0,
            'difference_qps': float(throughput2) - float(throughput1) if throughput1 and throughput2 else 0
        }

    def _compare_resources(self,
                          model1_info: Dict,
                          model2_info: Dict) -> Dict[str, float]:
        """Compare resource usage"""
        memory1 = model1_info.get('params', {}).get('memory_mb', 0)
        memory2 = model2_info.get('params', {}).get('memory_mb', 0)

        return {
            'model1_memory_mb': float(memory1) if memory1 else 0,
            'model2_memory_mb': float(memory2) if memory2 else 0,
            'difference_mb': float(memory2) - float(memory1) if memory1 and memory2 else 0
        }

    def _make_recommendation(self,
                            version1: str,
                            version2: str,
                            metric_comparisons: Dict,
                            statistically_significant: bool,
                            latency_comparison: Dict,
                            resource_comparison: Dict) -> Tuple[str, str, float]:
        """Make promotion recommendation"""

        reasons = []
        version2_score = 0
        version1_score = 0

        # Evaluate performance improvements
        for metric, comparison in metric_comparisons.items():
            improvement = comparison['improvement_percent']
            if improvement > 1:  # > 1% improvement
                version2_score += 1
                reasons.append(f"{metric} improved by {improvement:.2f}%")
            elif improvement < -1:  # > 1% degradation
                version1_score += 1
                reasons.append(f"{metric} degraded by {abs(improvement):.2f}%")

        # Consider statistical significance
        if not statistically_significant:
            reasons.append("Improvement not statistically significant")
            version2_score -= 2

        # Consider latency
        latency_diff = latency_comparison.get('difference_ms', 0)
        if latency_diff > 10:  # 10ms degradation
            version1_score += 1
            reasons.append(f"Latency increased by {latency_diff:.1f}ms")
        elif latency_diff < -10:  # 10ms improvement
            version2_score += 1
            reasons.append(f"Latency reduced by {abs(latency_diff):.1f}ms")

        # Make decision
        if version2_score > version1_score:
            recommended = version2
            confidence = min(version2_score / (version2_score + version1_score + 1), 0.99)
        elif version1_score > version2_score:
            recommended = version1
            confidence = min(version1_score / (version2_score + version1_score + 1), 0.99)
        else:
            recommended = version1  # Keep current if tie
            confidence = 0.5
            reasons.append("Models perform similarly - keeping current version")

        reason = "; ".join(reasons)
        return recommended, reason, confidence

    def _load_model(self, model_name: str, version: str):
        """Load model from registry"""
        model_uri = f"models:/{model_name}/{version}"
        return mlflow.sklearn.load_model(model_uri)

class PromotionManager:
    """
    Automated model promotion based on criteria

    Evaluates models against criteria and promotes automatically
    """

    def __init__(self,
                 comparator: ModelComparator,
                 governance: ModelGovernanceManager):
        """
        Initialize promotion manager

        Args:
            comparator: ModelComparator instance
            governance: ModelGovernanceManager instance
        """
        self.comparator = comparator
        self.governance = governance
        self.logger = logging.getLogger("PromotionManager")

    def evaluate_for_promotion(self,
                              model_name: str,
                              candidate_version: str,
                              current_version: str,
                              criteria: PromotionCriteria,
                              test_data: pd.DataFrame = None,
                              test_labels: np.ndarray = None) -> Dict:
        """
        Evaluate if candidate model should be promoted

        Args:
            model_name: Model name
            candidate_version: Candidate version
            current_version: Current production version
            criteria: Promotion criteria
            test_data: Test dataset
            test_labels: Test labels

        Returns:
            Evaluation results with promotion decision
        """
        self.logger.info(
            f"Evaluating {model_name} v{candidate_version} for promotion"
        )

        # Compare models
        comparison = self.comparator.compare_models(
            model_name=model_name,
            version1=current_version,
            version2=candidate_version,
            test_data=test_data,
            test_labels=test_labels
        )

        # Evaluate against criteria
        checks = {}
        all_passed = True

        # Check accuracy improvement
        accuracy_comp = comparison.metric_comparisons.get('accuracy', {})
        accuracy_improvement = accuracy_comp.get('improvement_percent', 0) / 100

        checks['accuracy_improvement'] = accuracy_improvement >= criteria.min_accuracy_improvement
        if not checks['accuracy_improvement']:
            all_passed = False
            self.logger.warning(
                f"Accuracy improvement {accuracy_improvement:.3f} below threshold "
                f"{criteria.min_accuracy_improvement}"
            )

        # Check F1 improvement
        f1_comp = comparison.metric_comparisons.get('f1_score', {})
        f1_improvement = f1_comp.get('improvement_percent', 0) / 100

        checks['f1_improvement'] = f1_improvement >= criteria.min_f1_improvement
        if not checks['f1_improvement']:
            all_passed = False

        # Check statistical significance
        if criteria.require_statistical_significance:
            checks['statistical_significance'] = comparison.statistically_significant
            if not checks['statistical_significance']:
                all_passed = False
                self.logger.warning("Improvement not statistically significant")

        # Check latency constraint
        if criteria.max_latency_ms:
            candidate_latency = comparison.latency_comparison.get('model2_ms', 0)
            checks['latency_constraint'] = candidate_latency <= criteria.max_latency_ms
            if not checks['latency_constraint']:
                all_passed = False

        # Check resource constraints
        if criteria.max_memory_mb:
            candidate_memory = comparison.resource_comparison.get('model2_memory_mb', 0)
            checks['memory_constraint'] = candidate_memory <= criteria.max_memory_mb
            if not checks['memory_constraint']:
                all_passed = False

        return {
            'should_promote': all_passed,
            'comparison': comparison,
            'criteria_checks': checks,
            'recommendation': comparison.recommendation_reason
        }

    def auto_promote(self,
                    model_name: str,
                    candidate_version: str,
                    current_version: str,
                    criteria: PromotionCriteria,
                    requester: str,
                    test_data: pd.DataFrame = None,
                    test_labels: np.ndarray = None) -> Optional[str]:
        """
        Automatically promote model if criteria met

        Args:
            model_name: Model name
            candidate_version: Candidate version
            current_version: Current version
            criteria: Promotion criteria
            requester: User requesting promotion
            test_data: Test data
            test_labels: Test labels

        Returns:
            Approval request ID if promotion initiated, None otherwise
        """
        # Evaluate for promotion
        evaluation = self.evaluate_for_promotion(
            model_name=model_name,
            candidate_version=candidate_version,
            current_version=current_version,
            criteria=criteria,
            test_data=test_data,
            test_labels=test_labels
        )

        if not evaluation['should_promote']:
            self.logger.info(
                f"Model {model_name} v{candidate_version} does not meet promotion criteria"
            )
            return None

        # Create approval request
        comparison = evaluation['comparison']

        request_id = self.governance.request_approval(
            model_name=model_name,
            version=candidate_version,
            source_stage='Staging',
            target_stage='Production',
            requester=requester,
            requester_role=UserRole.ML_ENGINEER,
            justification=f"Automated promotion: {comparison.recommendation_reason}",
            impact_assessment=f"Performance improvements: {comparison.metric_comparisons}",
            rollback_plan=f"Rollback to v{current_version} if issues detected"
        )

        self.logger.info(
            f"Created promotion request {request_id} for {model_name} v{candidate_version}"
        )

        return request_id

# Example usage
def example_model_comparison_and_promotion():
    """Example of model comparison and promotion"""

    # Initialize components
    registry = MLflowModelRegistry(tracking_uri="http://localhost:5000")
    comparator = ModelComparator(registry)
    governance = ModelGovernanceManager(registry)
    promotion_manager = PromotionManager(comparator, governance)

    # Create test data
    from sklearn.datasets import make_classification
    X_test, y_test = make_classification(
        n_samples=1000, n_features=20, random_state=42
    )

    # Compare models
    comparison = comparator.compare_models(
        model_name='fraud_detector',
        version1='2',  # Current production
        version2='3',  # Candidate
        test_data=pd.DataFrame(X_test),
        test_labels=y_test
    )

    print(f"Comparison Results:")
    print(f"  Recommended: {comparison.recommended_model}")
    print(f"  Reason: {comparison.recommendation_reason}")
    print(f"  Confidence: {comparison.confidence_score:.2f}")
    print(f"  Statistically significant: {comparison.statistically_significant}")

    # Define promotion criteria
    criteria = PromotionCriteria(
        min_accuracy_improvement=0.01,  # 1%
        min_f1_improvement=0.01,
        require_statistical_significance=True,
        max_latency_ms=100,
        max_memory_mb=2000
    )

    # Evaluate for promotion
    evaluation = promotion_manager.evaluate_for_promotion(
        model_name='fraud_detector',
        candidate_version='3',
        current_version='2',
        criteria=criteria,
        test_data=pd.DataFrame(X_test),
        test_labels=y_test
    )

    print(f"\nPromotion Evaluation:")
    print(f"  Should promote: {evaluation['should_promote']}")
    print(f"  Criteria checks: {evaluation['criteria_checks']}")

    # Auto-promote if criteria met
    if evaluation['should_promote']:
        request_id = promotion_manager.auto_promote(
            model_name='fraud_detector',
            candidate_version='3',
            current_version='2',
            criteria=criteria,
            requester='ml-pipeline@company.com',
            test_data=pd.DataFrame(X_test),
            test_labels=y_test
        )
        print(f"\nPromotion request created: {request_id}")
\end{lstlisting}

\section{CI/CD Pipeline Integration}

\subsection{Automated Model Lifecycle with CI/CD}

Integrating model registry with CI/CD pipelines automates model training, testing, registration, and deployment. This reduces manual errors, ensures reproducibility, and accelerates delivery.

\textbf{CI/CD integration benefits:}
\begin{itemize}
    \item \textbf{Automated testing:} Run validation tests on every model version
    \item \textbf{Consistent deployment:} Standardized promotion workflows
    \item \textbf{Version control:} Automatic model versioning and tagging
    \item \textbf{Audit trail:} Complete history of model changes
    \item \textbf{Rollback capability:} Quick reversion to previous versions
    \item \textbf{Infrastructure as code:} Declarative model deployment
\end{itemize}

\subsection{CI/CD Pipeline Implementation}

\begin{lstlisting}[language=Python, caption=CI/CD Pipeline Integration for ML Models]
from typing import Dict, List, Optional
from dataclasses import dataclass
import subprocess
import yaml
import json
import os
from pathlib import Path

@dataclass
class PipelineConfig:
    """CI/CD pipeline configuration"""
    # Repository
    git_repo: str
    git_branch: str

    # Training
    training_script: str
    training_data_path: str
    hyperparameters_file: str

    # Testing
    test_script: str
    test_data_path: str
    performance_thresholds: Dict[str, float]

    # Model registration
    model_name: str
    registry_uri: str

    # Deployment
    deployment_target: str  # staging, production
    kubernetes_namespace: str
    docker_image_registry: str

class MLPipeline:
    """
    ML CI/CD pipeline orchestrator

    Automates model training, testing, registration, and deployment
    """

    def __init__(self,
                 config: PipelineConfig,
                 registry: MLflowModelRegistry,
                 governance: ModelGovernanceManager):
        """
        Initialize ML pipeline

        Args:
            config: Pipeline configuration
            registry: MLflow registry
            governance: Governance manager
        """
        self.config = config
        self.registry = registry
        self.governance = governance
        self.logger = logging.getLogger("MLPipeline")

    def run_training_pipeline(self) -> Dict:
        """
        Run complete training pipeline

        Returns:
            Pipeline results
        """
        self.logger.info("Starting training pipeline")

        results = {
            'git_commit': None,
            'run_id': None,
            'model_version': None,
            'metrics': {},
            'tests_passed': False,
            'registered': False
        }

        try:
            # Step 1: Get git commit info
            results['git_commit'] = self._get_git_commit()

            # Step 2: Train model
            run_id, metrics = self._train_model()
            results['run_id'] = run_id
            results['metrics'] = metrics

            # Step 3: Run model tests
            test_results = self._run_model_tests(run_id)
            results['tests_passed'] = test_results['passed']

            if not test_results['passed']:
                self.logger.error("Model tests failed")
                return results

            # Step 4: Register model
            model_version = self._register_model(run_id, metrics)
            results['model_version'] = model_version
            results['registered'] = True

            # Step 5: Request approval for staging
            self._request_staging_approval(model_version)

            self.logger.info(
                f"Training pipeline completed successfully. Model v{model_version}"
            )

        except Exception as e:
            self.logger.error(f"Training pipeline failed: {str(e)}")
            raise

        return results

    def run_deployment_pipeline(self,
                                model_version: str,
                                target_stage: str) -> bool:
        """
        Run deployment pipeline

        Args:
            model_version: Model version to deploy
            target_stage: Target stage (Staging, Production)

        Returns:
            True if deployment successful
        """
        self.logger.info(
            f"Starting deployment pipeline for v{model_version} to {target_stage}"
        )

        try:
            # Step 1: Build Docker image
            image_tag = self._build_docker_image(model_version)

            # Step 2: Push to registry
            self._push_docker_image(image_tag)

            # Step 3: Deploy to Kubernetes
            self._deploy_to_kubernetes(image_tag, target_stage)

            # Step 4: Run smoke tests
            smoke_passed = self._run_smoke_tests(target_stage)

            if not smoke_passed:
                self.logger.error("Smoke tests failed, rolling back")
                self._rollback_deployment(target_stage)
                return False

            self.logger.info(f"Deployment to {target_stage} successful")
            return True

        except Exception as e:
            self.logger.error(f"Deployment failed: {str(e)}")
            self._rollback_deployment(target_stage)
            return False

    def _get_git_commit(self) -> str:
        """Get current git commit hash"""
        result = subprocess.run(
            ['git', 'rev-parse', 'HEAD'],
            capture_output=True,
            text=True
        )
        return result.stdout.strip()

    def _train_model(self) -> tuple:
        """Execute training script"""
        self.logger.info("Training model")

        # Load hyperparameters
        with open(self.config.hyperparameters_file, 'r') as f:
            hyperparameters = yaml.safe_load(f)

        # Run training script
        env = os.environ.copy()
        env['MLFLOW_TRACKING_URI'] = self.config.registry_uri

        result = subprocess.run(
            ['python', self.config.training_script,
             '--data', self.config.training_data_path,
             '--params', self.config.hyperparameters_file],
            env=env,
            capture_output=True,
            text=True
        )

        if result.returncode != 0:
            raise Exception(f"Training failed: {result.stderr}")

        # Parse output to get run_id and metrics
        output = json.loads(result.stdout)
        return output['run_id'], output['metrics']

    def _run_model_tests(self, run_id: str) -> Dict:
        """Run model validation tests"""
        self.logger.info("Running model tests")

        result = subprocess.run(
            ['python', self.config.test_script,
             '--run-id', run_id,
             '--test-data', self.config.test_data_path],
            capture_output=True,
            text=True
        )

        if result.returncode != 0:
            return {'passed': False, 'error': result.stderr}

        test_results = json.loads(result.stdout)

        # Check against thresholds
        passed = True
        for metric, threshold in self.config.performance_thresholds.items():
            actual = test_results['metrics'].get(metric, 0)
            if actual < threshold:
                self.logger.warning(
                    f"{metric}={actual:.3f} below threshold {threshold}"
                )
                passed = False

        return {'passed': passed, 'results': test_results}

    def _register_model(self, run_id: str, metrics: Dict) -> str:
        """Register model in MLflow"""
        self.logger.info("Registering model")

        # Load model from run
        model_uri = f"runs:/{run_id}/model"

        # Create metadata
        metadata = ModelMetadata(
            name=self.config.model_name,
            version="auto",
            stage="None",
            run_id=run_id,
            experiment_id="0",
            metrics=metrics,
            training_date=datetime.now(),
            training_duration_seconds=0,  # Would be tracked
            dataset_version="auto",
            dataset_size=0,
            model_type="auto",
            framework="sklearn",
            framework_version="1.2.0",
            python_version="3.9",
            dependencies=[],
            tags={'ci_cd': 'true', 'automated': 'true'},
            description=f"Automated build from commit {self._get_git_commit()[:8]}"
        )

        # Register
        version = self.registry.register_model(
            model=None,  # Already logged in training
            model_name=self.config.model_name,
            run_id=run_id,
            metadata=metadata
        )

        return version

    def _request_staging_approval(self, model_version: str):
        """Request approval for staging deployment"""
        self.governance.request_approval(
            model_name=self.config.model_name,
            version=model_version,
            source_stage='None',
            target_stage='Staging',
            requester='ci-cd-pipeline',
            requester_role=UserRole.ML_ENGINEER,
            justification='Automated build passed all tests',
            impact_assessment='New model version for staging validation',
            rollback_plan='Automated rollback on failure'
        )

    def _build_docker_image(self, model_version: str) -> str:
        """Build Docker image with model"""
        image_tag = f"{self.config.docker_image_registry}/{self.config.model_name}:{model_version}"

        self.logger.info(f"Building Docker image: {image_tag}")

        # Create Dockerfile
        dockerfile_content = f"""
FROM python:3.9-slim

# Install dependencies
COPY requirements.txt .
RUN pip install -r requirements.txt

# Copy model serving code
COPY serve.py .

# Set MLflow environment
ENV MLFLOW_TRACKING_URI={self.config.registry_uri}
ENV MODEL_NAME={self.config.model_name}
ENV MODEL_VERSION={model_version}

# Expose port
EXPOSE 8000

# Run server
CMD ["python", "serve.py"]
        """

        with open('Dockerfile', 'w') as f:
            f.write(dockerfile_content)

        # Build image
        result = subprocess.run(
            ['docker', 'build', '-t', image_tag, '.'],
            capture_output=True,
            text=True
        )

        if result.returncode != 0:
            raise Exception(f"Docker build failed: {result.stderr}")

        return image_tag

    def _push_docker_image(self, image_tag: str):
        """Push Docker image to registry"""
        self.logger.info(f"Pushing image: {image_tag}")

        result = subprocess.run(
            ['docker', 'push', image_tag],
            capture_output=True,
            text=True
        )

        if result.returncode != 0:
            raise Exception(f"Docker push failed: {result.stderr}")

    def _deploy_to_kubernetes(self, image_tag: str, target_stage: str):
        """Deploy to Kubernetes"""
        self.logger.info(f"Deploying to Kubernetes: {target_stage}")

        # Create Kubernetes deployment manifest
        deployment = {
            'apiVersion': 'apps/v1',
            'kind': 'Deployment',
            'metadata': {
                'name': f"{self.config.model_name}-{target_stage.lower()}",
                'namespace': self.config.kubernetes_namespace
            },
            'spec': {
                'replicas': 3 if target_stage == 'Production' else 1,
                'selector': {
                    'matchLabels': {
                        'app': self.config.model_name,
                        'stage': target_stage.lower()
                    }
                },
                'template': {
                    'metadata': {
                        'labels': {
                            'app': self.config.model_name,
                            'stage': target_stage.lower()
                        }
                    },
                    'spec': {
                        'containers': [{
                            'name': 'model-server',
                            'image': image_tag,
                            'ports': [{'containerPort': 8000}],
                            'resources': {
                                'requests': {
                                    'memory': '1Gi',
                                    'cpu': '500m'
                                },
                                'limits': {
                                    'memory': '2Gi',
                                    'cpu': '1000m'
                                }
                            }
                        }]
                    }
                }
            }
        }

        # Write manifest
        with open('deployment.yaml', 'w') as f:
            yaml.dump(deployment, f)

        # Apply to cluster
        result = subprocess.run(
            ['kubectl', 'apply', '-f', 'deployment.yaml'],
            capture_output=True,
            text=True
        )

        if result.returncode != 0:
            raise Exception(f"Kubernetes deployment failed: {result.stderr}")

    def _run_smoke_tests(self, target_stage: str) -> bool:
        """Run smoke tests on deployed model"""
        self.logger.info("Running smoke tests")

        # Would make test requests to deployed endpoint
        # Simplified here
        return True

    def _rollback_deployment(self, target_stage: str):
        """Rollback failed deployment"""
        self.logger.info(f"Rolling back {target_stage} deployment")

        # Rollback to previous version
        result = subprocess.run(
            ['kubectl', 'rollout', 'undo',
             f"deployment/{self.config.model_name}-{target_stage.lower()}",
             '-n', self.config.kubernetes_namespace],
            capture_output=True,
            text=True
        )

        if result.returncode != 0:
            self.logger.error(f"Rollback failed: {result.stderr}")

    def generate_github_actions_workflow(self) -> str:
        """Generate GitHub Actions workflow YAML"""
        workflow = {
            'name': 'ML Model CI/CD',
            'on': {
                'push': {
                    'branches': ['main', 'develop']
                },
                'pull_request': {
                    'branches': ['main']
                }
            },
            'jobs': {
                'train-and-test': {
                    'runs-on': 'ubuntu-latest',
                    'steps': [
                        {'uses': 'actions/checkout@v2'},
                        {
                            'name': 'Set up Python',
                            'uses': 'actions/setup-python@v2',
                            'with': {'python-version': '3.9'}
                        },
                        {
                            'name': 'Install dependencies',
                            'run': 'pip install -r requirements.txt'
                        },
                        {
                            'name': 'Train model',
                            'env': {
                                'MLFLOW_TRACKING_URI': '${{ secrets.MLFLOW_TRACKING_URI }}'
                            },
                            'run': f'python {self.config.training_script}'
                        },
                        {
                            'name': 'Run tests',
                            'run': f'python {self.config.test_script}'
                        },
                        {
                            'name': 'Register model',
                            'if': "github.ref == 'refs/heads/main'",
                            'run': 'python scripts/register_model.py'
                        }
                    ]
                },
                'deploy-staging': {
                    'needs': 'train-and-test',
                    'if': "github.ref == 'refs/heads/develop'",
                    'runs-on': 'ubuntu-latest',
                    'steps': [
                        {
                            'name': 'Deploy to Staging',
                            'run': 'python scripts/deploy.py --stage staging'
                        }
                    ]
                },
                'deploy-production': {
                    'needs': 'train-and-test',
                    'if': "github.ref == 'refs/heads/main'",
                    'runs-on': 'ubuntu-latest',
                    'steps': [
                        {
                            'name': 'Deploy to Production',
                            'run': 'python scripts/deploy.py --stage production'
                        }
                    ]
                }
            }
        }

        return yaml.dump(workflow, default_flow_style=False)

# Example usage
def example_cicd_pipeline():
    """Example of CI/CD pipeline integration"""

    # Configuration
    config = PipelineConfig(
        git_repo='https://github.com/company/ml-models',
        git_branch='main',
        training_script='train.py',
        training_data_path='data/train.csv',
        hyperparameters_file='config/hyperparameters.yaml',
        test_script='test_model.py',
        test_data_path='data/test.csv',
        performance_thresholds={
            'accuracy': 0.90,
            'f1_score': 0.85,
            'auc': 0.88
        },
        model_name='fraud_detector',
        registry_uri='http://mlflow.company.com',
        deployment_target='production',
        kubernetes_namespace='ml-models',
        docker_image_registry='gcr.io/company'
    )

    # Initialize pipeline
    registry = MLflowModelRegistry(tracking_uri=config.registry_uri)
    governance = ModelGovernanceManager(registry)
    pipeline = MLPipeline(config, registry, governance)

    # Run training pipeline
    results = pipeline.run_training_pipeline()

    if results['registered']:
        print(f"Model registered: v{results['model_version']}")
        print(f"Metrics: {results['metrics']}")

        # Deploy to staging
        if results['tests_passed']:
            success = pipeline.run_deployment_pipeline(
                model_version=results['model_version'],
                target_stage='Staging'
            )
            print(f"Staging deployment: {'Success' if success else 'Failed'}")

    # Generate GitHub Actions workflow
    workflow_yaml = pipeline.generate_github_actions_workflow()
    print("\nGenerated GitHub Actions workflow:")
    print(workflow_yaml)
\end{lstlisting}

\section{Model Packaging and Containerization}

\subsection{Reproducible Model Packaging}

Proper model packaging ensures reproducibility across environments. Packaging includes the model artifacts, dependencies, runtime configuration, and metadata needed for deployment.

\textbf{Packaging requirements:}
\begin{itemize}
    \item \textbf{Model artifacts:} Serialized model files (pickle, ONNX, SavedModel)
    \item \textbf{Dependencies:} Exact versions of all libraries
    \item \textbf{Environment:} Python version, system libraries
    \item \textbf{Preprocessing:} Feature transformations, encoders, scalers
    \item \textbf{Configuration:} Model hyperparameters, thresholds
    \item \textbf{Metadata:} Model signature, lineage, performance metrics
\end{itemize}

\subsection{Model Packaging and Container Implementation}

\begin{lstlisting}[language=Python, caption=Model Packaging and Containerization System]
from typing import Dict, List, Optional, Any
from pathlib import Path
import pickle
import json
import shutil
from dataclasses import dataclass

@dataclass
class PackageConfig:
    """Model package configuration"""
    model_name: str
    model_version: str
    model_format: str  # mlflow, onnx, pickle, savedmodel
    include_preprocessing: bool = True
    include_postprocessing: bool = True
    optimize_size: bool = True
    compression: str = "gzip"  # gzip, zip, tar

class ModelPackager:
    """
    Package models in multiple formats with dependencies

    Creates self-contained model packages for deployment
    """

    def __init__(self, output_dir: Path = Path("./model_packages")):
        """
        Initialize model packager

        Args:
            output_dir: Directory for packaged models
        """
        self.output_dir = output_dir
        self.output_dir.mkdir(parents=True, exist_ok=True)
        self.logger = logging.getLogger("ModelPackager")

    def package_model(self,
                     model: Any,
                     config: PackageConfig,
                     metadata: Dict,
                     preprocessing_pipeline: Any = None,
                     postprocessing_fn: Optional[callable] = None) -> Path:
        """
        Package model with all dependencies

        Args:
            model: Trained model
            config: Package configuration
            metadata: Model metadata
            preprocessing_pipeline: Sklearn pipeline or custom preprocessor
            postprocessing_fn: Post-prediction processing function

        Returns:
            Path to packaged model
        """
        self.logger.info(f"Packaging {config.model_name} v{config.model_version}")

        # Create package directory
        package_dir = self.output_dir / f"{config.model_name}_{config.model_version}"
        package_dir.mkdir(exist_ok=True)

        # Save model in specified format
        if config.model_format == 'mlflow':
            self._save_mlflow_format(model, package_dir, metadata)
        elif config.model_format == 'onnx':
            self._save_onnx_format(model, package_dir, metadata)
        elif config.model_format == 'pickle':
            self._save_pickle_format(model, package_dir)
        else:
            raise ValueError(f"Unsupported format: {config.model_format}")

        # Save preprocessing pipeline
        if config.include_preprocessing and preprocessing_pipeline:
            self._save_preprocessing(preprocessing_pipeline, package_dir)

        # Save postprocessing
        if config.include_postprocessing and postprocessing_fn:
            self._save_postprocessing(postprocessing_fn, package_dir)

        # Save metadata
        self._save_metadata(metadata, package_dir)

        # Generate requirements.txt
        self._generate_requirements(package_dir)

        # Create inference script
        self._create_inference_script(package_dir, config)

        # Compress if requested
        if config.optimize_size:
            archive_path = self._compress_package(package_dir, config.compression)
            return archive_path

        return package_dir

    def _save_mlflow_format(self, model: Any, package_dir: Path, metadata: Dict):
        """Save in MLflow format"""
        import mlflow.sklearn

        mlflow.sklearn.save_model(
            model,
            path=str(package_dir / "model"),
            signature=metadata.get('signature')
        )

    def _save_onnx_format(self, model: Any, package_dir: Path, metadata: Dict):
        """Convert and save in ONNX format"""
        try:
            from skl2onnx import convert_sklearn
            from skl2onnx.common.data_types import FloatTensorType

            # Infer input shape from metadata
            n_features = metadata.get('n_features', 10)
            initial_type = [('float_input', FloatTensorType([None, n_features]))]

            # Convert to ONNX
            onnx_model = convert_sklearn(model, initial_types=initial_type)

            # Save
            with open(package_dir / "model.onnx", "wb") as f:
                f.write(onnx_model.SerializeToString())

            self.logger.info("Model converted to ONNX format")

        except ImportError:
            self.logger.warning("skl2onnx not installed, falling back to pickle")
            self._save_pickle_format(model, package_dir)

    def _save_pickle_format(self, model: Any, package_dir: Path):
        """Save in pickle format"""
        with open(package_dir / "model.pkl", "wb") as f:
            pickle.dump(model, f)

    def _save_preprocessing(self, pipeline: Any, package_dir: Path):
        """Save preprocessing pipeline"""
        with open(package_dir / "preprocessing.pkl", "wb") as f:
            pickle.dump(pipeline, f)

    def _save_postprocessing(self, fn: callable, package_dir: Path):
        """Save postprocessing function"""
        import dill

        with open(package_dir / "postprocessing.pkl", "wb") as f:
            dill.dump(fn, f)

    def _save_metadata(self, metadata: Dict, package_dir: Path):
        """Save model metadata"""
        metadata_file = package_dir / "metadata.json"

        with open(metadata_file, "w") as f:
            json.dump(metadata, f, indent=2, default=str)

    def _generate_requirements(self, package_dir: Path):
        """Generate requirements.txt"""
        import pkg_resources

        # Get installed packages
        installed_packages = {pkg.key: pkg.version for pkg in pkg_resources.working_set}

        # Core ML packages
        required_packages = [
            'numpy', 'pandas', 'scikit-learn', 'mlflow',
            'scipy', 'joblib', 'fastapi', 'uvicorn'
        ]

        requirements = []
        for package in required_packages:
            if package in installed_packages:
                requirements.append(f"{package}=={installed_packages[package]}")

        # Write requirements
        with open(package_dir / "requirements.txt", "w") as f:
            f.write("\n".join(requirements))

    def _create_inference_script(self, package_dir: Path, config: PackageConfig):
        """Create standalone inference script"""
        script_content = f"""
import pickle
import json
import numpy as np
from pathlib import Path

class ModelInference:
    def __init__(self, model_dir: str):
        self.model_dir = Path(model_dir)

        # Load model
        if (self.model_dir / 'model.pkl').exists():
            with open(self.model_dir / 'model.pkl', 'rb') as f:
                self.model = pickle.load(f)
        else:
            import mlflow.sklearn
            self.model = mlflow.sklearn.load_model(str(self.model_dir / 'model'))

        # Load preprocessing if exists
        preproc_path = self.model_dir / 'preprocessing.pkl'
        self.preprocessing = None
        if preproc_path.exists():
            with open(preproc_path, 'rb') as f:
                self.preprocessing = pickle.load(f)

        # Load metadata
        with open(self.model_dir / 'metadata.json', 'r') as f:
            self.metadata = json.load(f)

    def predict(self, input_data):
        # Preprocess
        if self.preprocessing:
            input_data = self.preprocessing.transform(input_data)

        # Predict
        predictions = self.model.predict(input_data)

        return predictions

if __name__ == '__main__':
    # Example usage
    inference = ModelInference('.')

    # Make prediction
    import pandas as pd
    test_data = pd.DataFrame(np.random.rand(1, 10))
    predictions = inference.predict(test_data)
    print(f"Predictions: {{predictions}}")
"""

        with open(package_dir / "inference.py", "w") as f:
            f.write(script_content)

    def _compress_package(self, package_dir: Path, compression: str) -> Path:
        """Compress package directory"""
        archive_name = f"{package_dir.name}.tar.gz"
        archive_path = self.output_dir / archive_name

        shutil.make_archive(
            str(package_dir),
            'gztar',
            self.output_dir,
            package_dir.name
        )

        self.logger.info(f"Package compressed: {archive_path}")
        return archive_path

class ContainerBuilder:
    """
    Build optimized containers for model deployment

    Creates production-ready Docker images
    """

    def __init__(self):
        """Initialize container builder"""
        self.logger = logging.getLogger("ContainerBuilder")

    def build_optimized_container(self,
                                  model_package_path: Path,
                                  model_name: str,
                                  version: str,
                                  base_image: str = "python:3.9-slim",
                                  optimize: bool = True) -> str:
        """
        Build optimized Docker container

        Args:
            model_package_path: Path to model package
            model_name: Model name
            version: Model version
            base_image: Base Docker image
            optimize: Apply optimization techniques

        Returns:
            Image tag
        """
        self.logger.info(f"Building container for {model_name}:{version}")

        # Generate optimized Dockerfile
        dockerfile = self._generate_dockerfile(
            model_package_path,
            base_image,
            optimize
        )

        # Write Dockerfile
        with open('Dockerfile', 'w') as f:
            f.write(dockerfile)

        # Build image
        image_tag = f"{model_name}:{version}"
        result = subprocess.run(
            ['docker', 'build', '-t', image_tag, '.'],
            capture_output=True,
            text=True
        )

        if result.returncode != 0:
            raise Exception(f"Docker build failed: {result.stderr}")

        self.logger.info(f"Container built: {image_tag}")
        return image_tag

    def _generate_dockerfile(self,
                            model_package_path: Path,
                            base_image: str,
                            optimize: bool) -> str:
        """Generate optimized Dockerfile"""

        # Optimization layers
        optimization_layers = ""
        if optimize:
            optimization_layers = """
# Optimization: Use multi-stage build
FROM python:3.9 as builder
WORKDIR /build
COPY requirements.txt .
RUN pip install --user -r requirements.txt

# Production stage
"""

        dockerfile = f"""
{optimization_layers}
FROM {base_image}

# Set working directory
WORKDIR /app

# Copy dependencies from builder (if optimized)
{"COPY --from=builder /root/.local /root/.local" if optimize else ""}

# Install dependencies
COPY requirements.txt .
{"" if optimize else "RUN pip install --no-cache-dir -r requirements.txt"}

# Copy model package
COPY {model_package_path} ./model/

# Copy serving code
COPY inference.py .
COPY serve.py .

# Create non-root user for security
RUN useradd -m -u 1000 modeluser && chown -R modeluser:modeluser /app
USER modeluser

# Set environment variables
ENV MODEL_DIR=/app/model
ENV PYTHONPATH=/app

# Health check
HEALTHCHECK --interval=30s --timeout=3s --start-period=5s --retries=3 \\
  CMD python -c "import requests; requests.get('http://localhost:8000/health')"

# Expose port
EXPOSE 8000

# Run server
CMD ["python", "serve.py"]
"""

        return dockerfile

# Example usage
def example_packaging_and_containerization():
    """Example of model packaging and containerization"""

    from sklearn.ensemble import RandomForestClassifier
    from sklearn.preprocessing import StandardScaler
    from sklearn.pipeline import Pipeline

    # Train a simple model
    X_train = np.random.rand(100, 10)
    y_train = np.random.randint(0, 2, 100)

    # Create preprocessing pipeline
    preprocessing = StandardScaler()
    preprocessing.fit(X_train)

    # Train model
    model = RandomForestClassifier(n_estimators=100, random_state=42)
    model.fit(preprocessing.transform(X_train), y_train)

    # Package configuration
    config = PackageConfig(
        model_name='fraud_detector',
        model_version='v1.0.0',
        model_format='mlflow',
        include_preprocessing=True,
        optimize_size=True,
        compression='gzip'
    )

    # Metadata
    metadata = {
        'model_type': 'RandomForestClassifier',
        'framework': 'scikit-learn',
        'n_features': 10,
        'metrics': {
            'accuracy': 0.95,
            'f1_score': 0.93
        },
        'created_at': datetime.now().isoformat()
    }

    # Package model
    packager = ModelPackager()
    package_path = packager.package_model(
        model=model,
        config=config,
        metadata=metadata,
        preprocessing_pipeline=preprocessing
    )

    print(f"Model packaged at: {package_path}")

    # Build container
    builder = ContainerBuilder()
    image_tag = builder.build_optimized_container(
        model_package_path=package_path,
        model_name=config.model_name,
        version=config.model_version,
        optimize=True
    )

    print(f"Container built: {image_tag}")
\end{lstlisting}

\section{Exercises}

\subsection{Exercise 1: Model Registration and Lifecycle Management}

\textbf{Problem:} Implement a complete model registration workflow that trains a model, registers it in MLflow, transitions through stages (Staging → Production), and maintains version history.

\textbf{Solution:}
\begin{lstlisting}[language=Python]
# Exercise 1 Solution
from sklearn.datasets import load_breast_cancer
from sklearn.model_selection import train_test_split
from sklearn.ensemble import GradientBoostingClassifier
from sklearn.metrics import accuracy_score, f1_score, roc_auc_score
import mlflow
import mlflow.sklearn

def solution_exercise_1():
    """Complete model registration and lifecycle workflow"""

    # Initialize MLflow
    mlflow.set_tracking_uri("http://localhost:5000")
    mlflow.set_experiment("breast_cancer_detection")

    # Load data
    data = load_breast_cancer()
    X_train, X_test, y_train, y_test = train_test_split(
        data.data, data.target, test_size=0.2, random_state=42
    )

    # Train model
    model = GradientBoostingClassifier(
        n_estimators=100,
        learning_rate=0.1,
        max_depth=3,
        random_state=42
    )

    with mlflow.start_run(run_name="breast_cancer_v1") as run:
        # Train
        model.fit(X_train, y_train)

        # Evaluate
        y_pred = model.predict(X_test)
        y_pred_proba = model.predict_proba(X_test)[:, 1]

        accuracy = accuracy_score(y_test, y_pred)
        f1 = f1_score(y_test, y_pred)
        auc = roc_auc_score(y_test, y_pred_proba)

        # Log metrics
        mlflow.log_metric("accuracy", accuracy)
        mlflow.log_metric("f1_score", f1)
        mlflow.log_metric("auc_roc", auc)

        # Log parameters
        mlflow.log_param("n_estimators", 100)
        mlflow.log_param("learning_rate", 0.1)
        mlflow.log_param("max_depth", 3)

        # Log model
        mlflow.sklearn.log_model(
            model,
            artifact_path="model",
            registered_model_name="breast_cancer_detector"
        )

        run_id = run.info.run_id

    # Initialize registry manager
    registry = MLflowModelRegistry(tracking_uri="http://localhost:5000")

    # Get latest version
    versions = registry.client.get_latest_versions("breast_cancer_detector")
    latest_version = versions[0].version

    # Transition to Staging
    registry.transition_model_stage(
        model_name="breast_cancer_detector",
        version=latest_version,
        stage="Staging"
    )

    print(f"Model v{latest_version} moved to Staging")

    # After validation in staging, move to Production
    registry.transition_model_stage(
        model_name="breast_cancer_detector",
        version=latest_version,
        stage="Production",
        archive_existing=True
    )

    print(f"Model v{latest_version} promoted to Production")

    # Get model stages summary
    summary = registry.get_model_stages_summary("breast_cancer_detector")
    print(f"Model stages: {summary}")

solution_exercise_1()
\end{lstlisting}

\subsection{Exercise 2: Complete Model Lineage Tracking}

\textbf{Problem:} Implement comprehensive lineage tracking for a model that captures data provenance, code version, training configuration, and dependencies.

\textbf{Solution:}
\begin{lstlisting}[language=Python]
# Exercise 2 Solution
import pandas as pd
import numpy as np
from sklearn.ensemble import RandomForestClassifier

def solution_exercise_2():
    """Implement complete lineage tracking"""

    # Create training dataset
    np.random.seed(42)
    X_train = pd.DataFrame({
        'feature_1': np.random.rand(1000),
        'feature_2': np.random.rand(1000),
        'feature_3': np.random.randint(0, 10, 1000),
        'feature_4': np.random.rand(1000),
        'feature_5': np.random.randint(0, 5, 1000)
    })
    y_train = np.random.randint(0, 2, 1000)

    # Initialize lineage tracker
    tracker = LineageTracker(storage_backend="mlflow")

    # Track data lineage
    data_lineage = tracker.track_data_lineage(
        dataset=X_train,
        dataset_name="customer_churn_features",
        dataset_version="v2024-01-20",
        data_source="warehouse.customer_events",
        preprocessing_steps=[
            "remove_null_values",
            "normalize_numerical_features",
            "encode_categorical_features",
            "handle_class_imbalance_smote",
            "train_test_split_80_20"
        ]
    )

    print(f"Data lineage tracked:")
    print(f"  Dataset: {data_lineage.dataset_name}")
    print(f"  Version: {data_lineage.dataset_version}")
    print(f"  Hash: {data_lineage.dataset_hash[:16]}...")
    print(f"  Samples: {data_lineage.num_samples}")
    print(f"  Features: {data_lineage.num_features}")
    print(f"  Preprocessing: {data_lineage.preprocessing_steps}")

    # Track code lineage
    code_lineage = tracker.track_code_lineage(
        training_script_paths=["train.py", "preprocessing.py", "utils.py"]
    )

    print(f"\nCode lineage tracked:")
    print(f"  Git commit: {code_lineage.git_commit[:8]}")
    print(f"  Git branch: {code_lineage.git_branch}")
    print(f"  Repository: {code_lineage.git_repo}")
    print(f"  Dependencies: {len(code_lineage.dependencies)} packages")

    # Train model
    model = RandomForestClassifier(
        n_estimators=200,
        max_depth=15,
        min_samples_split=10,
        random_state=42
    )
    model.fit(X_train, y_train)

    # Define hyperparameters
    hyperparameters = {
        'n_estimators': 200,
        'max_depth': 15,
        'min_samples_split': 10,
        'random_state': 42,
        'criterion': 'gini',
        'bootstrap': True
    }

    # Calculate metrics
    from sklearn.metrics import accuracy_score, precision_score, recall_score
    y_pred = model.predict(X_train)  # Simplified - should use validation set

    metrics = {
        'accuracy': accuracy_score(y_train, y_pred),
        'precision': precision_score(y_train, y_pred),
        'recall': recall_score(y_train, y_pred)
    }

    # Create complete lineage
    model_lineage = tracker.create_lineage(
        model_id="churn_model_abc123",
        model_name="customer_churn_predictor",
        model_version="v3.0",
        data_lineage=data_lineage,
        code_lineage=code_lineage,
        hyperparameters=hyperparameters,
        training_duration=245.5,
        metrics=metrics,
        parent_models=[]  # No parent models (not transfer learning)
    )

    print(f"\nComplete model lineage created:")
    print(f"  Model ID: {model_lineage.model_id}")
    print(f"  Model: {model_lineage.model_name} v{model_lineage.model_version}")
    print(f"  Hardware: {model_lineage.hardware}")
    print(f"  Training duration: {model_lineage.training_duration}s")

    return model_lineage

lineage = solution_exercise_2()
\end{lstlisting}

\subsection{Exercise 3: Model Governance with Approval Workflows}

\textbf{Problem:} Implement a governance system with validation gates and multi-stage approval workflow for promoting a model to production.

\textbf{Solution:}
\begin{lstlisting}[language=Python]
# Exercise 3 Solution

def solution_exercise_3():
    """Implement model governance with approval workflows"""

    # Initialize components
    registry = MLflowModelRegistry(tracking_uri="http://localhost:5000")
    governance = ModelGovernanceManager(registry)

    # Define custom validation gates
    def validate_minimum_samples(model_info: Dict, request: ApprovalRequest) -> bool:
        """Ensure model trained on sufficient data"""
        dataset_size = model_info.get('params', {}).get('dataset_size', 0)
        min_required = 10000

        if int(dataset_size) < min_required:
            return False
        return True

    def validate_model_bias(model_info: Dict, request: ApprovalRequest) -> bool:
        """Check for model bias metrics"""
        tags = model_info.get('tags', {})

        # Check if bias assessment was performed
        if tags.get('bias_assessed') != 'true':
            return False

        # Check fairness metrics
        demographic_parity = float(tags.get('demographic_parity', '0'))
        if abs(demographic_parity) > 0.1:  # Max 10% disparity
            return False

        return True

    def validate_security_scan(model_info: Dict, request: ApprovalRequest) -> bool:
        """Validate security scanning completed"""
        tags = model_info.get('tags', {})
        return tags.get('security_scanned') == 'true'

    # Register validation gates for Production
    governance.register_validation_gate(
        target_stage='Production',
        gate=ValidationGate(
            name='performance_check',
            description='Validate performance meets thresholds',
            validator_func=validate_model_performance,
            required=True
        )
    )

    governance.register_validation_gate(
        target_stage='Production',
        gate=ValidationGate(
            name='sample_size_check',
            description='Ensure sufficient training data',
            validator_func=validate_minimum_samples,
            required=True
        )
    )

    governance.register_validation_gate(
        target_stage='Production',
        gate=ValidationGate(
            name='bias_check',
            description='Check for model bias and fairness',
            validator_func=validate_model_bias,
            required=True
        )
    )

    governance.register_validation_gate(
        target_stage='Production',
        gate=ValidationGate(
            name='security_check',
            description='Validate security scan completion',
            validator_func=validate_security_scan,
            required=False  # Optional for now
        )
    )

    # Request approval for production
    request_id = governance.request_approval(
        model_name='fraud_detector',
        version='5',
        source_stage='Staging',
        target_stage='Production',
        requester='data-scientist@company.com',
        requester_role=UserRole.DATA_SCIENTIST,
        justification='New model achieves 3% improvement in precision with no recall degradation',
        impact_assessment='Expected to reduce false positives by 20%, saving $500K annually in manual review costs',
        rollback_plan='Automated rollback to v4 if production accuracy drops below 0.92 or latency exceeds 50ms'
    )

    print(f"Approval request created: {request_id}")

    # Check request status
    status = governance.get_request_status(request_id)
    print(f"\nRequest Status:")
    print(f"  Status: {status['status']}")
    print(f"  Model: {status['model_name']} v{status['version']}")
    print(f"  Target: {status['target_stage']}")
    print(f"  Validation results: {status['validation_results']}")

    # Team lead approval
    governance.approve_request(
        request_id=request_id,
        approver='team-lead@company.com',
        approver_role=UserRole.TEAM_LEAD,
        comments='Validated model performance in staging. Metrics look good. Approved.'
    )

    print(f"\nTeam Lead approved")

    # Production owner approval
    governance.approve_request(
        request_id=request_id,
        approver='prod-owner@company.com',
        approver_role=UserRole.PRODUCTION_OWNER,
        comments='Infrastructure capacity confirmed. Load testing passed. Final approval granted.'
    )

    print(f"Production Owner approved")

    # Check final status
    final_status = governance.get_request_status(request_id)
    print(f"\nFinal Status: {final_status['status']}")
    print(f"Approvals: {len(final_status['approvals'])}")

    # Get audit trail
    audit_trail = governance.get_audit_trail(model_name='fraud_detector')
    print(f"\nAudit Trail ({len(audit_trail)} events):")
    for event in audit_trail[-5:]:  # Show last 5 events
        print(f"  [{event['timestamp']}] {event['event_type']} by {event['user']}")

    return request_id

request_id = solution_exercise_3()
\end{lstlisting}

\subsection{Exercise 4: Statistical Model Comparison}

\textbf{Problem:} Implement rigorous statistical comparison between two model versions, including McNemar's test and bootstrap confidence intervals.

\textbf{Solution:}
\begin{lstlisting}[language=Python]
# Exercise 4 Solution
from sklearn.datasets import make_classification
from sklearn.ensemble import RandomForestClassifier, GradientBoostingClassifier

def solution_exercise_4():
    """Statistical model comparison with hypothesis testing"""

    # Create test data
    X_test, y_test = make_classification(
        n_samples=2000,
        n_features=20,
        n_informative=15,
        n_redundant=5,
        random_state=42
    )

    # Train two models
    model_v1 = RandomForestClassifier(n_estimators=100, random_state=42)
    model_v2 = GradientBoostingClassifier(n_estimators=100, random_state=42)

    # Simplified training (use actual train set in practice)
    model_v1.fit(X_test[:1000], y_test[:1000])
    model_v2.fit(X_test[:1000], y_test[:1000])

    # Initialize comparator
    registry = MLflowModelRegistry(tracking_uri="http://localhost:5000")
    comparator = ModelComparator(registry)

    # Register both models first (simplified here)
    # In practice, load from registry

    # Create comparison with test data
    print("Comparing Model v1 (RandomForest) vs Model v2 (GradientBoosting)")
    print("=" * 70)

    # Get predictions
    pred_v1_proba = model_v1.predict_proba(X_test)[:, 1]
    pred_v2_proba = model_v2.predict_proba(X_test)[:, 1]

    pred_v1 = (pred_v1_proba > 0.5).astype(int)
    pred_v2 = (pred_v2_proba > 0.5).astype(int)

    # Calculate metrics
    from sklearn.metrics import accuracy_score, precision_score, recall_score, f1_score, roc_auc_score

    metrics_v1 = {
        'accuracy': accuracy_score(y_test, pred_v1),
        'precision': precision_score(y_test, pred_v1),
        'recall': recall_score(y_test, pred_v1),
        'f1_score': f1_score(y_test, pred_v1),
        'auc_roc': roc_auc_score(y_test, pred_v1_proba)
    }

    metrics_v2 = {
        'accuracy': accuracy_score(y_test, pred_v2),
        'precision': precision_score(y_test, pred_v2),
        'recall': recall_score(y_test, pred_v2),
        'f1_score': f1_score(y_test, pred_v2),
        'auc_roc': roc_auc_score(y_test, pred_v2_proba)
    }

    print("\nPerformance Metrics:")
    print(f"{'Metric':<15} {'v1 (RF)':<12} {'v2 (GB)':<12} {'Improvement':<15}")
    print("-" * 70)

    for metric in metrics_v1.keys():
        v1_val = metrics_v1[metric]
        v2_val = metrics_v2[metric]
        improvement = ((v2_val - v1_val) / v1_val) * 100

        print(f"{metric:<15} {v1_val:<12.4f} {v2_val:<12.4f} {improvement:>+.2f}%")

    # McNemar's Test
    print("\n\nStatistical Significance Testing:")
    print("-" * 70)

    # Contingency table
    both_correct = np.sum((pred_v1 == y_test) & (pred_v2 == y_test))
    only_v1_correct = np.sum((pred_v1 == y_test) & (pred_v2 != y_test))
    only_v2_correct = np.sum((pred_v1 != y_test) & (pred_v2 == y_test))
    both_wrong = np.sum((pred_v1 != y_test) & (pred_v2 != y_test))

    print(f"\nContingency Table:")
    print(f"  Both models correct: {both_correct}")
    print(f"  Only v1 correct: {only_v1_correct}")
    print(f"  Only v2 correct: {only_v2_correct}")
    print(f"  Both models wrong: {both_wrong}")

    # McNemar statistic
    if only_v1_correct + only_v2_correct > 0:
        from scipy import stats
        mcnemar_stat = ((abs(only_v1_correct - only_v2_correct) - 1) ** 2) / \
                      (only_v1_correct + only_v2_correct)
        mcnemar_p = 1 - stats.chi2.cdf(mcnemar_stat, 1)

        print(f"\nMcNemar's Test:")
        print(f"  Chi-square statistic: {mcnemar_stat:.4f}")
        print(f"  p-value: {mcnemar_p:.4f}")
        print(f"  Significant (p < 0.05): {mcnemar_p < 0.05}")

    # Bootstrap confidence intervals
    print(f"\nBootstrap Confidence Intervals (95%):")

    n_bootstrap = 1000
    auc_diffs = []

    for _ in range(n_bootstrap):
        indices = np.random.choice(len(y_test), len(y_test), replace=True)
        boot_auc1 = roc_auc_score(y_test[indices], pred_v1_proba[indices])
        boot_auc2 = roc_auc_score(y_test[indices], pred_v2_proba[indices])
        auc_diffs.append(boot_auc2 - boot_auc1)

    auc_diffs = np.array(auc_diffs)
    ci_lower = np.percentile(auc_diffs, 2.5)
    ci_upper = np.percentile(auc_diffs, 97.5)

    print(f"  AUC difference: ({ci_lower:.4f}, {ci_upper:.4f})")
    print(f"  CI excludes zero: {ci_lower > 0 or ci_upper < 0}")

    # Recommendation
    print(f"\n\nRecommendation:")
    if mcnemar_p < 0.05 and ci_lower > 0:
        print(f"  [OK] Model v2 shows statistically significant improvement")
        print(f"  [OK] Recommend promoting v2 to production")
    elif mcnemar_p >= 0.05:
        print(f"  [X] No statistically significant difference detected")
        print(f"  [X] Keep current model v1 in production")
    else:
        print(f"  ? Mixed results - require additional validation")

solution_exercise_4()
\end{lstlisting}

\subsection{Exercise 5: CI/CD Pipeline for Model Deployment}

\textbf{Problem:} Create a complete CI/CD pipeline that trains a model, runs tests, registers it, and deploys to Kubernetes.

\textbf{Solution:}
\begin{lstlisting}[language=Python]
# Exercise 5 Solution

def solution_exercise_5():
    """Complete CI/CD pipeline implementation"""

    # Pipeline configuration
    config = PipelineConfig(
        git_repo='https://github.com/company/ml-fraud-detection',
        git_branch='main',
        training_script='scripts/train_fraud_detector.py',
        training_data_path='data/fraud_training.csv',
        hyperparameters_file='config/fraud_detector_params.yaml',
        test_script='scripts/test_fraud_detector.py',
        test_data_path='data/fraud_test.csv',
        performance_thresholds={
            'accuracy': 0.93,
            'precision': 0.90,
            'recall': 0.88,
            'f1_score': 0.89,
            'auc_roc': 0.92
        },
        model_name='fraud_detector',
        registry_uri='http://mlflow.company.com',
        deployment_target='production',
        kubernetes_namespace='ml-models-prod',
        docker_image_registry='gcr.io/company-ml'
    )

    # Initialize pipeline components
    registry = MLflowModelRegistry(tracking_uri=config.registry_uri)
    governance = ModelGovernanceManager(registry)
    pipeline = MLPipeline(config, registry, governance)

    print("=" * 70)
    print("ML CI/CD Pipeline Execution")
    print("=" * 70)

    # Step 1: Run training pipeline
    print("\n[1/3] Running Training Pipeline...")
    print("-" * 70)

    try:
        results = pipeline.run_training_pipeline()

        if not results['tests_passed']:
            print("[X] Training pipeline failed: Tests did not pass")
            print(f"  Metrics: {results['metrics']}")
            return False

        print(f"[OK] Training completed successfully")
        print(f"  Git commit: {results['git_commit'][:8]}")
        print(f"  Run ID: {results['run_id']}")
        print(f"  Model version: {results['model_version']}")
        print(f"  Metrics: {results['metrics']}")

    except Exception as e:
        print(f"[X] Training pipeline failed: {str(e)}")
        return False

    # Step 2: Run deployment to staging
    print(f"\n[2/3] Deploying to Staging...")
    print("-" * 70)

    try:
        staging_success = pipeline.run_deployment_pipeline(
            model_version=results['model_version'],
            target_stage='Staging'
        )

        if not staging_success:
            print("[X] Staging deployment failed")
            return False

        print(f"[OK] Deployed to Staging successfully")
        print(f"  Model: {config.model_name}:{results['model_version']}")
        print(f"  Namespace: {config.kubernetes_namespace}")

    except Exception as e:
        print(f"[X] Staging deployment failed: {str(e)}")
        return False

    # Step 3: Production promotion (requires approval)
    print(f"\n[3/3] Production Promotion Workflow...")
    print("-" * 70)

    # Request production approval
    request_id = governance.request_approval(
        model_name=config.model_name,
        version=results['model_version'],
        source_stage='Staging',
        target_stage='Production',
        requester='ci-cd-pipeline@company.com',
        requester_role=UserRole.ML_ENGINEER,
        justification=f"Automated deployment - all tests passed with metrics: {results['metrics']}",
        impact_assessment="Improved fraud detection accuracy expected to save $2M annually",
        rollback_plan="Automated rollback to previous version if production metrics degrade"
    )

    print(f"[OK] Production approval request created: {request_id}")
    print(f"  Awaiting approvals from: Team Lead, Production Owner")
    print(f"\n  Next steps:")
    print(f"    1. Team Lead reviews and approves")
    print(f"    2. Production Owner reviews and approves")
    print(f"    3. Model automatically promoted to Production")
    print(f"    4. Kubernetes deployment updated")

    # Generate GitHub Actions workflow
    print(f"\n\nGenerated GitHub Actions Workflow:")
    print("=" * 70)
    workflow = pipeline.generate_github_actions_workflow()
    print(workflow[:500] + "...")  # Show first 500 chars

    return True

success = solution_exercise_5()
print(f"\n\nPipeline execution: {'SUCCESS' if success else 'FAILED'}")
\end{lstlisting}

\subsection{Exercise 6: Model Packaging for Multi-Format Deployment}

\textbf{Problem:} Package a model in multiple formats (MLflow, ONNX, pickle) with preprocessing pipelines and create an optimized Docker container.

\textbf{Solution:}
\begin{lstlisting}[language=Python]
# Exercise 6 Solution
from sklearn.preprocessing import StandardScaler, OneHotEncoder
from sklearn.compose import ColumnTransformer

def solution_exercise_6():
    """Package model in multiple formats with containerization"""

    # Create sample model with preprocessing
    X_train = pd.DataFrame({
        'age': np.random.randint(18, 80, 1000),
        'income': np.random.randint(20000, 200000, 1000),
        'credit_score': np.random.randint(300, 850, 1000),
        'category': np.random.choice(['A', 'B', 'C'], 1000)
    })
    y_train = np.random.randint(0, 2, 1000)

    # Create preprocessing pipeline
    numeric_features = ['age', 'income', 'credit_score']
    categorical_features = ['category']

    preprocessor = ColumnTransformer(
        transformers=[
            ('num', StandardScaler(), numeric_features),
            ('cat', OneHotEncoder(drop='first'), categorical_features)
        ]
    )

    # Fit preprocessing
    preprocessor.fit(X_train)
    X_processed = preprocessor.transform(X_train)

    # Train model
    model = RandomForestClassifier(
        n_estimators=150,
        max_depth=12,
        min_samples_split=8,
        random_state=42
    )
    model.fit(X_processed, y_train)

    # Package in multiple formats
    packager = ModelPackager(output_dir=Path("./model_packages"))

    formats = ['mlflow', 'pickle', 'onnx']
    packages = {}

    print("Packaging model in multiple formats...")
    print("=" * 70)

    for fmt in formats:
        print(f"\nPackaging in {fmt.upper()} format...")

        config = PackageConfig(
            model_name='credit_risk_model',
            model_version='v2.1.0',
            model_format=fmt,
            include_preprocessing=True,
            include_postprocessing=False,
            optimize_size=True,
            compression='gzip'
        )

        metadata = {
            'model_type': 'RandomForestClassifier',
            'framework': 'scikit-learn',
            'framework_version': '1.2.0',
            'n_features': len(numeric_features) + len(categorical_features),
            'feature_names': X_train.columns.tolist(),
            'numeric_features': numeric_features,
            'categorical_features': categorical_features,
            'metrics': {
                'accuracy': 0.94,
                'precision': 0.92,
                'recall': 0.91,
                'f1_score': 0.91
            },
            'created_at': datetime.now().isoformat(),
            'created_by': 'automated-pipeline',
            'use_case': 'credit_risk_assessment'
        }

        package_path = packager.package_model(
            model=model,
            config=config,
            metadata=metadata,
            preprocessing_pipeline=preprocessor
        )

        packages[fmt] = package_path
        print(f"  [OK] Package created: {package_path}")

        # Check package contents
        if package_path.is_dir():
            contents = list(package_path.glob('*'))
            print(f"  Package contents:")
            for item in contents:
                print(f"    - {item.name}")

    # Build Docker containers for each format
    print(f"\n\nBuilding Docker Containers...")
    print("=" * 70)

    builder = ContainerBuilder()

    for fmt, package_path in packages.items():
        if package_path.suffix == '.gz':
            # Extract archive first
            import tarfile
            extract_dir = package_path.parent / package_path.stem.replace('.tar', '')
            with tarfile.open(package_path, 'r:gz') as tar:
                tar.extractall(path=extract_dir)
            package_path = extract_dir

        print(f"\nBuilding container for {fmt.upper()} format...")

        image_tag = builder.build_optimized_container(
            model_package_path=package_path,
            model_name=f'credit_risk_model_{fmt}',
            version='v2.1.0',
            base_image='python:3.9-slim',
            optimize=True
        )

        print(f"  [OK] Container built: {image_tag}")
        print(f"  Ready for deployment to Kubernetes/Cloud Run")

    # Summary
    print(f"\n\nPackaging Summary:")
    print("=" * 70)
    print(f"Model: credit_risk_model v2.1.0")
    print(f"Formats packaged: {len(packages)}")
    print(f"Containers built: {len(packages)}")
    print(f"\nDeployment options:")
    print(f"  - MLflow format: Best for MLflow-based serving")
    print(f"  - ONNX format: Cross-platform, optimized inference")
    print(f"  - Pickle format: Simple Python-based serving")

    return packages

packages = solution_exercise_6()
\end{lstlisting}

\section{Chapter Summary}

This chapter covered comprehensive model versioning and registry systems essential for production ML operations. We explored the complete lifecycle of model management from registration to deployment, focusing on practical implementations that ensure reproducibility, governance, and operational excellence.

\subsection{Key Concepts Covered}

\textbf{MLflow Model Registry:}
\begin{itemize}
    \item Model registration with comprehensive metadata and signatures
    \item Stage-based lifecycle management (None, Staging, Production, Archived)
    \item Version control and model comparison capabilities
    \item Model loading and serving from registry
\end{itemize}

\textbf{Model Lineage and Provenance:}
\begin{itemize}
    \item Data lineage tracking with dataset versioning and hashing
    \item Code lineage capturing git commits, branches, and dependencies
    \item Complete model provenance including hyperparameters and training metadata
    \item Reproducibility through comprehensive artifact tracking
\end{itemize}

\textbf{Governance and Approval Workflows:}
\begin{itemize}
    \item Multi-stage approval processes with role-based access control
    \item Automated validation gates ensuring model quality
    \item Audit trails for compliance and regulatory requirements
    \item Integration with ticketing systems for formal approvals
\end{itemize}

\textbf{Model Comparison and Promotion:}
\begin{itemize}
    \item Statistical significance testing with McNemar's test
    \item Bootstrap confidence intervals for metric comparisons
    \item Comprehensive evaluation across performance, latency, and resource usage
    \item Automated promotion based on criteria and thresholds
\end{itemize}

\textbf{CI/CD Integration:}
\begin{itemize}
    \item Automated training, testing, and registration pipelines
    \item Docker containerization and Kubernetes deployment
    \item GitHub Actions and Jenkins workflow integration
    \item Automated rollback and smoke testing
\end{itemize}

\textbf{Model Packaging and Containerization:}
\begin{itemize}
    \item Multi-format packaging (MLflow, ONNX, pickle)
    \item Dependency management and requirements generation
    \item Optimized Docker images with multi-stage builds
    \item Security best practices with non-root users
\end{itemize}

\subsection{Best Practices}

\begin{enumerate}
    \item \textbf{Always track complete lineage:} Capture data, code, and training provenance for every model version to ensure reproducibility
    \item \textbf{Implement validation gates:} Automated checks prevent low-quality models from reaching production
    \item \textbf{Require statistical significance:} Use hypothesis testing to validate that model improvements are real, not random variation
    \item \textbf{Automate everything:} CI/CD pipelines reduce manual errors and accelerate deployment
    \item \textbf{Maintain audit trails:} Complete history of model changes satisfies compliance requirements
    \item \textbf{Package with preprocessing:} Include all transformation steps in model packages for deployment consistency
    \item \textbf{Use staged rollouts:} Test models in staging before production deployment
    \item \textbf{Monitor production models:} Continuous validation ensures models maintain performance
\end{enumerate}

\subsection{Common Pitfalls to Avoid}

\begin{itemize}
    \item \textbf{Incomplete metadata:} Missing training details make reproduction impossible
    \item \textbf{Skipping approval gates:} Direct production deployments bypass critical validation
    \item \textbf{Ignoring statistical significance:} Small metric improvements may not be meaningful
    \item \textbf{Manual deployments:} Error-prone and slow compared to automated pipelines
    \item \textbf{Missing rollback plans:} Production issues without rollback capability cause outages
    \item \textbf{Inadequate testing:} Insufficient validation leads to production failures
    \item \textbf{Poor version control:} Unclear model history makes debugging difficult
\end{itemize}

\subsection{Production Checklist}

Before deploying a model registry system to production, verify:

\begin{enumerate}
    \item Model registry infrastructure is highly available and backed up
    \item All model versions have complete metadata and lineage
    \item Approval workflows are configured with appropriate roles
    \item Validation gates are implemented and tested
    \item CI/CD pipelines are operational with proper error handling
    \item Docker images are scanned for security vulnerabilities
    \item Kubernetes deployments have resource limits and health checks
    \item Rollback procedures are documented and tested
    \item Monitoring and alerting are configured for model performance
    \item Audit logging captures all model lifecycle events
\end{enumerate}

\subsection{Next Steps}

The model versioning and registry system forms the foundation for production ML operations. In subsequent chapters, we'll explore:

\begin{itemize}
    \item Model monitoring and observability systems
    \item Data drift detection and model retraining triggers
    \item Multi-model ensemble management
    \item Cost optimization for model serving
    \item Advanced deployment strategies (shadow, canary, A/B)
\end{itemize}

By implementing the practices and systems covered in this chapter, teams can confidently manage dozens or hundreds of model versions, ensure reproducibility, maintain governance, and deploy models safely to production. The combination of MLflow, automated workflows, and statistical rigor provides a production-grade foundation for ML operations at scale.
